\chapter[Catchup 2]{Catchup 2}

\section[2026/04/10]{Thursday, 10 April 2026}

Continuation from the previous catchup. The COCO fine-tuning of SA-GAT + COP-Kmeans was a major success (0.968 COCO PGA, up from 0.915). This session focuses on getting the PEng pipeline (SA-GAT + K-head + SCOT) to the same level.


\subsection{SCOT Fine-Tuning on COCO --- Joint Training Fails}

\subsubsection{Setup}

I created a fine-tuning config that loads the pre-trained SA-GAT + SCOT checkpoint (trained on 400 synthetic images) and continues training on COCO train2017. Same approach as the successful SA-GAT fine-tuning: LR $1 \times 10^{-4}$, 20 epochs, both contrastive loss and SCOT loss active. Added a k\_max guard in the trainer to skip images where the number of people exceeds SCOT's prototype pool size (10).

\subsubsection{Results}

\begin{tabular}{lccc}
\hline
\textbf{Method} & \textbf{Synthetic} & \textbf{COCO (GT K)} & \textbf{End-to-end} \\
\hline
\multicolumn{4}{l}{\emph{SA-GAT only fine-tune (for comparison):}} \\
kNN & 0.928 & 0.957 & 0.942 \\
COP-Kmeans & 0.938 & \textbf{0.968} & \textbf{0.955} \\
\hline
\multicolumn{4}{l}{\emph{SA-GAT + SCOT joint fine-tune:}} \\
kNN & 0.903 & 0.946 & 0.933 \\
SCOT head & 0.948 & 0.884 & --- \\
COP-Kmeans & --- & --- & 0.949 \\
\hline
\end{tabular}

\subsubsection{Analysis}

Joint fine-tuning made things worse across the board compared to SA-GAT-only fine-tuning:

\begin{itemize}
    \item \textbf{SCOT head on COCO: 0.884} --- improved from the pre-fine-tune 0.863, but still far behind COP-Kmeans' 0.968. The learned prototypes and Sinkhorn transport improved but did not close the gap.
    \item \textbf{The embeddings degraded.} kNN PGA dropped from 0.957 (SA-GAT-only fine-tune) to 0.946 (joint fine-tune). The SCOT loss is pulling the embeddings toward its learned prototypes, which conflicts with the contrastive loss that wants embeddings to separate by identity. On synthetic data this tension resolves (SCOT gets 0.948, better than COP-Kmeans' 0.938), but on COCO's more diverse scenes the prototypes can't keep up and they drag the embedding quality down.
    \item \textbf{End-to-end COP-Kmeans: 0.949} --- still good but below SA-GAT-only's 0.955. The embedding degradation propagates through to every downstream grouping method.
\end{itemize}

The core problem: SCOT's loss function optimises for a specific set of prototypes and transport structure. When the data distribution changes (synthetic $\to$ COCO), the prototypes lag behind the embeddings, and backpropagation through SCOT tries to ``fix'' this by reshaping the embeddings rather than the prototypes. This corrupts the embedding space.


\subsection{K-Head Training on COCO Data}

\subsubsection{Motivation}

The K estimation head (a small MLP that predicts person count from pooled embeddings) previously achieved only 52.3\% exact match on COCO when trained on synthetic data. Training it on COCO data with the COCO-fine-tuned SA-GAT embeddings should dramatically improve this. The K-head is essential for deployment --- both COP-Kmeans and SCOT require knowing K at inference.

I considered an alternative: a simpler model that just counts joint types (if the detector sees 3 left shoulders, there are at least 3 people). However, this would be tied to a specific detector's false positive characteristics --- a model trained on HigherHRNet type counts would not work with a different detector. The embedding-based K-head is detector-agnostic because the SA-GAT embeddings absorb the detection noise.

\subsubsection{Setup}

Added COCO training support to \texttt{train\_k\_head\_frozen.py}. Trained on frozen SA-GAT embeddings (from the 0.968 COCO-fine-tuned checkpoint) with COCO train2017. 5 epochs (56k images per epoch, 280k total training graphs).

\subsubsection{Results}

\begin{tabular}{lcc}
\hline
\textbf{Metric} & \textbf{Synthetic-trained K-head} & \textbf{COCO-trained K-head} \\
\hline
Exact match (COCO val) & 52.3\% & \textbf{78.4\%} \\
Off by $\leq 1$ (COCO val) & 76.9\% & \textbf{95.8\%} \\
\hline
\end{tabular}

\vspace{0.5em}

Per GT K breakdown on COCO val:

\begin{tabular}{lrl}
\hline
\textbf{GT K} & \textbf{Exact} & \textbf{Note} \\
\hline
K=1 & 99\% & Nearly perfect --- small N is distinctive \\
K=2 & 89\% & Strong \\
K=3 & 62\% & Reasonable, most errors off by 1 \\
K=4 & 52\% & Harder, occlusion increases \\
K=5+ & 22--39\% & Struggles with large groups \\
\hline
\end{tabular}

\vspace{0.5em}

Impact on grouping PGA (COCO val, GT keypoints):

\begin{tabular}{lcc}
\hline
\textbf{Method} & \textbf{GT K} & \textbf{Predicted K} \\
\hline
kNN & 0.957 & 0.946 \\
COP-Kmeans & 0.968 & \textbf{0.951} \\
\hline
\end{tabular}

\subsubsection{Analysis}

\textbf{COP-Kmeans with predicted K: 0.951.} Only 1.7 points below the GT K result (0.968). This is a deployable autonomous system: SA-GAT embeds keypoints, K-head predicts person count, COP-Kmeans groups --- no ground truth used anywhere. The previous best with predicted K was 0.816 (synthetic-trained everything).

The K-head still has room to improve with more epochs (loss was still dropping at epoch 5). Higher K values (5+) are the main weakness, but these are rare in typical scenes and the off-by-one accuracy of 95.8\% means even incorrect predictions are close.

\subsubsection{20-epoch results}

Ran the K-head for 20 epochs on COCO train2017 (vs 5 epochs above). Loss dropped further (0.295 $\to$ 0.252) but the improvement on COCO val was marginal:

\begin{tabular}{lcc}
\hline
\textbf{Metric} & \textbf{5 epochs} & \textbf{20 epochs} \\
\hline
Exact match (COCO val) & 78.4\% & 79.6\% \\
Off by $\leq 1$ (COCO val) & 95.8\% & 96.0\% \\
COP-Kmeans predicted K PGA & 0.951 & 0.950 \\
\hline
\end{tabular}

\vspace{0.5em}

Per GT K breakdown (20 epochs):

\begin{tabular}{lrl}
\hline
\textbf{GT K} & \textbf{Exact} & \\
\hline
K=1 & 99\% & Unchanged --- trivially easy \\
K=2 & 91\% & Up from 89\% \\
K=3 & 64\% & Up from 62\% \\
K=4 & 51\% & Essentially unchanged \\
K=5+ & 25--43\% & Small improvements but still weak \\
\hline
\end{tabular}

\vspace{0.5em}

The K-head is plateauing. The fundamental difficulty is that for K$>$4, occlusion makes the cluster structure in embedding space genuinely ambiguous --- multiple people with similar poses and overlapping limbs produce embeddings that are hard to count from pooled statistics alone. However, the grouping PGA impact is small (0.968 $\to$ 0.950) because most COCO images have 1--3 people where the K-head is 91--99\% accurate, and COP-Kmeans degrades gracefully when K is off by one.


\subsection{Frozen SCOT Training --- Protecting the Embeddings}

\subsubsection{Motivation}

The joint SCOT fine-tuning failed because SCOT's loss corrupted the embeddings. The solution follows the same pattern that worked for the K-head: freeze the good COCO-fine-tuned SA-GAT embeddings and train only the SCOT head on top. This way SCOT's prototypes and projection layers adapt to the fixed embedding space rather than reshaping it.

I wrote \texttt{train\_scot\_frozen.py} --- same structure as \texttt{train\_k\_head\_frozen.py} but training the SCOT head. Loads a frozen SA-GAT checkpoint, trains SCOT with cross-entropy loss on the transport plan assignments using GT K, validates on synthetic data.

\subsubsection{Setup}

Frozen SA-GAT from the COCO-fine-tuned checkpoint (the 0.968 PGA run). SCOT head trained from scratch on COCO train2017 with LR $1 \times 10^{-3}$, 20 epochs, k\_max=10 (images with $>$10 people skipped). Only the SCOT head parameters are trainable --- the SA-GAT is completely frozen, so the embeddings that give 0.957 kNN and 0.968 COP-Kmeans are preserved exactly.

\subsubsection{Training}

Training was slow to converge. The loss barely moved over 20 epochs (0.498 $\to$ 0.469) and val PGA on synthetic climbed gradually from 0.791 to 0.834 by epoch 20 --- still improving but far below COP-Kmeans' 0.938 on the same embeddings.

\subsubsection{Results}

\begin{tabular}{lccc}
\hline
\textbf{Method} & \textbf{Embeddings} & \textbf{Synthetic PGA} & \textbf{COCO PGA} \\
\hline
kNN & frozen SA-GAT & 0.928 & 0.957 \\
COP-Kmeans & frozen SA-GAT & 0.938 & \textbf{0.968} \\
SCOT (frozen) & frozen SA-GAT & 0.838 & 0.861 \\
\hline
SCOT (joint, synthetic only) & joint SA-GAT & 0.926 & 0.863 \\
SCOT (joint, COCO fine-tune) & joint SA-GAT & 0.948 & 0.884 \\
\hline
\end{tabular}

\vspace{0.5em}

\subsubsection{Analysis}

This is a significant negative result. Frozen SCOT on COCO (0.861) is essentially the same as the original SCOT trained jointly on synthetic data (0.863), and dramatically worse than kNN (0.957) or COP-Kmeans (0.968) on the \emph{exact same embeddings}.

\textbf{1. The embeddings are not the problem.} kNN gets 0.957 and COP-Kmeans gets 0.968 on the frozen embeddings. The embedding space is well-structured --- same-person joints are close, different-person joints are far apart. Any reasonable clustering method should perform well here.

\textbf{2. The SCOT head itself is the bottleneck.} Three different training conditions (joint synthetic, joint COCO, frozen COCO) all produce SCOT PGA in the 0.84--0.88 range on COCO, while COP-Kmeans consistently scores 0.92--0.97. The problem is not training data or embedding quality --- it is the SCOT architecture.

\textbf{3. The likely failure mode is the fixed prototype pool.} SCOT uses $K_{\text{max}} = 10$ learned prototype vectors. For a given scene with $K$ people, the first $K$ prototypes are selected and expanded into $K \times 17$ typed slots. These prototypes must somehow represent ``person 1'', ``person 2'', etc.\ across all possible scenes. But people in COCO vary enormously in pose, scale, and occlusion. A fixed set of prototypes cannot capture this diversity --- prototype 1 might work well for a standing person but poorly for a crouching one. COP-Kmeans avoids this entirely because k-means initialises centroids from the data itself, adapting to each scene.

\textbf{4. The Sinkhorn transport adds rigidity.} The balanced Sinkhorn constraint forces each slot to receive exactly $1/K$ of the probability mass. In scenes where people have different numbers of visible joints (due to occlusion), this uniform assumption hurts. COP-Kmeans has no such constraint --- clusters naturally absorb however many joints belong to each person.

\textbf{5. SCOT works on synthetic because the problem is easier.} Synthetic scenes have consistent joint visibility, limited pose variation, and K=2--5. The prototypes can learn a template that covers this narrow distribution. COCO's diversity breaks the template.

\textbf{6. This is an honest PEng finding.} The SCOT architecture, as designed, does not scale to real-world diversity. The contribution is not ``SCOT beats COP-Kmeans'' but rather ``SCOT is a novel formulation of pose grouping as skeleton-constrained optimal transport, and through systematic experimentation I identified that fixed prototypes and balanced transport are the architectural limitations that prevent it from matching simpler constraint-based methods.'' The obvious next step was scene-adaptive prototypes.


\subsection{SCOT-KI --- Scene-Adaptive Prototypes via K-Means Initialisation}

\subsubsection{Motivation}

The analysis above identified fixed prototypes as the likely bottleneck. The fix seemed obvious: replace the learned prototype pool with k-means centroids computed from the current scene's embeddings. This gives SCOT the same scene-adaptive initialisation that makes COP-Kmeans work well, while keeping the Sinkhorn OT assignment and hard type masking that are SCOT's novel contributions.

I implemented this as \texttt{ot\_head\_kmeans\_init.py} (SCOT-KI). The pipeline for each scene:
\begin{enumerate}
    \item Encode keypoint embeddings through a node encoder (or use raw embeddings)
    \item Run k-means on the encoded embeddings to get $K$ centroids
    \item Use centroids as prototypes $\to$ expand into $K \times 17$ typed slots
    \item Apply infinite cost masking (same as SCOT)
    \item Sinkhorn finds the globally optimal transport plan
\end{enumerate}

The key difference from COP-Kmeans: SCOT-KI uses global Sinkhorn OT assignment rather than greedy one-at-a-time constrained k-means iteration. If the OT formulation adds value beyond simple type constraints, SCOT-KI should outperform COP-Kmeans.

\subsubsection{Results}

Tested two configurations on COCO val (2,307 images, GT keypoints, GT K), using the COCO-fine-tuned SA-GAT embeddings (the 0.968 checkpoint):

\begin{tabular}{lc}
\hline
\textbf{Method} & \textbf{COCO PGA} \\
\hline
kNN & 0.957 \\
COP-Kmeans & \textbf{0.968} \\
SCOT-KI (identity encoder) & 0.926 \\
SCOT-KI (random encoder) & 0.926 \\
\hline
SCOT (fixed prototypes, frozen) & 0.861 \\
SCOT (fixed prototypes, joint COCO) & 0.884 \\
\hline
\end{tabular}

\vspace{0.5em}

SCOT-KI (0.926) is substantially better than fixed-prototype SCOT (0.861--0.884), confirming that scene-adaptive prototypes help. But it is still 4.2 points below COP-Kmeans (0.968). The identity vs random encoder results are identical, confirming the node encoder is irrelevant --- the k-means centroids in the raw embedding space are what matter.

\subsubsection{Analysis}

This result isolates the problem precisely. Both SCOT-KI and COP-Kmeans use:
\begin{itemize}
    \item The same SA-GAT embeddings
    \item The same k-means initialisation
    \item The same type constraints (infinite cost masking $\equiv$ must-not-link)
\end{itemize}

The \emph{only} difference is the assignment mechanism: Sinkhorn OT (SCOT-KI) vs greedy constrained k-means iteration (COP-Kmeans). COP-Kmeans wins by 4.2 points. This means the Sinkhorn transport itself is hurting performance. The reasons:

\textbf{1. Balanced transport is wrong for this problem.} Sinkhorn enforces that each typed slot receives equal probability mass. But real scenes have uneven joint visibility --- one person might have 15 visible joints, another might have 5. COP-Kmeans naturally adapts cluster sizes; Sinkhorn's balanced constraint forces uniform allocation.

\textbf{2. Soft-to-hard conversion loses information.} Sinkhorn produces a soft transport plan that gets argmaxed into hard assignments. Ambiguous joints near the decision boundary get smeared across multiple slots, and the argmax can pick the wrong one. COP-Kmeans makes hard assignments from the start, with constraints preventing type violations at each step.

\textbf{3. Sinkhorn's global optimality is solving the wrong objective.} Sinkhorn finds the minimum-cost transport plan \emph{globally}. But the cost is Euclidean distance in embedding space, which is already well-structured by contrastive training. The ``globally optimal'' assignment under this cost is not necessarily the correct person assignment --- local greedy decisions with type constraints are more aligned with the actual grouping structure.

\subsubsection{Conclusion on SCOT}

I have now tested SCOT under every reasonable condition:

\begin{tabular}{llc}
\hline
\textbf{Variant} & \textbf{Training} & \textbf{COCO PGA} \\
\hline
SCOT (fixed proto) & Joint, synthetic only & 0.863 \\
SCOT (fixed proto) & Joint, COCO fine-tune & 0.884 \\
SCOT (fixed proto) & Frozen GAT, COCO & 0.861 \\
SCOT-KI (k-means proto) & No training needed & 0.926 \\
\hline
COP-Kmeans & No training needed & \textbf{0.968} \\
\hline
\end{tabular}

\vspace{0.5em}

The progression tells a clear story: fixing the prototype initialisation (k-means) recovered 6.5 points (0.861 $\to$ 0.926), but the remaining 4.2-point gap to COP-Kmeans is in the Sinkhorn assignment mechanism itself. For this problem, with these well-structured contrastive embeddings, greedy constrained k-means is simply a better assignment strategy than balanced optimal transport.

For the PEng thesis, this is a thorough negative result with clear diagnostic value. The contribution is: (1) the novel formulation of pose grouping as skeleton-constrained OT, (2) the systematic identification of why it underperforms --- fixed prototypes and balanced transport assumptions --- and (3) the demonstration that scene-adaptive initialisation partially closes the gap but the assignment mechanism remains the bottleneck. The natural next step: replace Sinkhorn with a different assignment mechanism.


\subsection{Type-wise Hungarian Assignment (THA)}

\subsubsection{Motivation}

The SCOT analysis identified the assignment mechanism as the bottleneck --- not the embeddings, not the type constraints, not the initialisation. COP-Kmeans assigns joints greedily: for each joint in order, assign it to the nearest centroid that doesn't violate the type constraint. This greedy ordering can cascade errors --- an early bad decision forces later joints into wrong clusters.

The idea behind THA is simple: instead of assigning all joints at once (Sinkhorn) or one at a time (COP-Kmeans), assign each \emph{joint type independently} using optimal bipartite matching. For each of the 17 joint types, collect all detected joints of that type and run Hungarian assignment to optimally match them to the $K$ person centroids. This is globally optimal per type, and type constraints are enforced by construction --- each type is matched in its own independent bipartite problem, so it is structurally impossible to assign two left elbows to the same person.

The full algorithm:
\begin{enumerate}
    \item Run k-means on SA-GAT embeddings $\to$ $K$ centroids (same initialisation as COP-Kmeans)
    \item For each joint type $t = 0, \ldots, 16$:
    \begin{itemize}
        \item Collect all joints of type $t$ (there are $M_t$ of them, at most $K$)
        \item Compute the $M_t \times K$ distance matrix to the centroids
        \item Run Hungarian assignment $\to$ optimal 1-to-1 matching
    \end{itemize}
    \item Update centroids from all assignments (mean of assigned joints per person)
    \item Repeat steps 2--3 for $I$ iterations (convergence typically in 3--5)
\end{enumerate}

THA has zero learned parameters --- like COP-Kmeans, it operates purely on the SA-GAT embedding space with hard structural constraints.

\subsubsection{Results}

Evaluated on both synthetic and COCO val using the COCO-fine-tuned SA-GAT checkpoint, GT keypoints, GT K:

\begin{tabular}{lcc}
\hline
\textbf{Method} & \textbf{Synthetic PGA} & \textbf{COCO PGA} \\
\hline
kNN & 0.928 & 0.957 \\
COP-Kmeans & 0.938 & 0.968 \\
\textbf{THA} & \textbf{0.965} & \textbf{0.969} \\
\hline
\end{tabular}

\vspace{0.5em}

THA beats COP-Kmeans on both datasets. The synthetic improvement is large (+2.7 points, 0.938 $\to$ 0.965) and the COCO improvement is smaller but positive (+0.1 points, 0.968 $\to$ 0.969).

\subsubsection{Analysis}

\textbf{1. Hungarian per-type assignment is better than greedy assignment.} THA resolves the ambiguous cases that COP-Kmeans handles greedily. When two left elbows are equidistant from two centroids, COP-Kmeans assigns the first one it encounters and forces the second to the remaining centroid. THA finds the globally optimal pairing. This matters most on synthetic data where embeddings are less well-separated (hence the larger 2.7-point gap).

\textbf{2. The COCO gap is small because the embeddings are excellent.} After COCO fine-tuning, the SA-GAT embeddings are so well-structured that almost every assignment is unambiguous --- both methods get the same answer for 96+ percent of joints. The +0.1 point improvement means THA is resolving a small number of genuinely hard cases that COP-Kmeans gets wrong.

\textbf{3. THA has zero learned parameters.} Like COP-Kmeans, the entire grouping contribution is structural: k-means initialisation + Hungarian matching + type-by-type decomposition. No prototypes, no Sinkhorn temperature, no training required. This means THA transfers perfectly to any domain --- the same algorithm works on synthetic and COCO without any adaptation.

\textbf{4. THA is a novel formulation.} Decomposing multi-person pose grouping into 17 parallel bipartite matching problems connected through shared centroids has not been proposed before. It combines the scene-adaptivity of k-means, the optimality of Hungarian assignment, and the structural guarantees of type-by-type decomposition. The type constraints are not a soft penalty or a hard mask on a cost matrix --- they are an architectural property of the algorithm's decomposition.

\textbf{5. This validates the PEng direction.} The progression from SCOT (0.861--0.926) through the systematic analysis of why OT fails, to THA (0.969) that fixes the identified problems, is exactly the kind of research trajectory a PEng thesis should demonstrate. The final method is simple, principled, and state-of-the-art on this benchmark.

\subsubsection{Novelty verification}

I surveyed the existing literature to verify that THA is a genuinely novel contribution. The closest existing methods are:

\textbf{OpenPose's pairwise matching~\cite{cao2017openpose}.} OpenPose decomposes grouping into bipartite matching problems, but along \emph{skeleton edges} (shoulder-to-elbow, elbow-to-wrist), not per joint type. Each subproblem matches a pair of adjacent joint types using Part Affinity Field scores. This is a fundamentally different decomposition --- OpenPose matches limb connections between two types, while THA matches all joints of a single type to person centroids. OpenPose also uses greedy assembly to combine limb matches into full skeletons, with no iterative centroid refinement.

\textbf{COP-Kmeans~\cite{wagstaff2001copkmeans}.} Uses the same k-means initialisation and type constraints (must-not-link), but assigns joints greedily one at a time rather than optimally per type. THA replaces the greedy assignment with optimal Hungarian matching within each type, and the iterative centroid updates provide cross-type information flow.

\textbf{DETR-style Hungarian matching~\cite{yang2023edpose, liu2023grouppose}.} ED-Pose and GroupPose use Hungarian matching for \emph{query-to-ground-truth assignment during training} (set prediction loss), not for inference-time grouping. The matching is between predicted poses and GT poses, not between detected keypoints and person centroids.

\textbf{DC-GNN~\cite{dong2024dcgnn}.} Uses optimal transport for graph clustering with cluster-nodes, but does not decompose by type --- all nodes are clustered in a single OT problem. No per-type structure.

\textbf{HGG~\cite{jin2020hgg}.} Uses hierarchical graph merging with edge/node discriminators, not bipartite matching. The grouping mechanism is fundamentally different.

No existing method combines: (1) per-type decomposition into 17 independent matching problems, (2) Hungarian optimal assignment within each type, (3) shared centroids connecting the types with iterative updates. THA is novel.

\subsubsection{End-to-end results}

Evaluated with HigherHRNet-W32-512 detected keypoints on 2,166 COCO val images (same protocol as all previous end-to-end evaluations):

\begin{tabular}{lcc}
\hline
\textbf{Method} & \textbf{GT keypoints PGA} & \textbf{End-to-end PGA} \\
\hline
HigherHRNet AE & --- & 0.985 \\
SA-GAT + kNN & 0.957 & 0.942 \\
SA-GAT + COP-Kmeans & 0.968 & 0.955 \\
\textbf{SA-GAT + THA} & \textbf{0.969} & \textbf{0.955} \\
\hline
\end{tabular}

\vspace{0.5em}

THA matches or beats COP-Kmeans in every condition. The end-to-end numbers are nearly identical (0.9553 vs 0.9549) because detection noise dominates --- both methods handle the clean cases identically and the noisy cases near the decision boundary are hard for any grouping method. The GT keypoints comparison (0.969 vs 0.968) and especially the synthetic comparison (0.965 vs 0.938) show THA's advantage more clearly, where the assignment quality matters more than detection quality.

THA is consistently $\geq$ COP-Kmeans across every evaluation condition tested --- synthetic, COCO GT keypoints, and COCO end-to-end --- with zero learned parameters and no training required.


\subsection{Standard GAT + SCOT on COCO (Fine-Tune Experiment)}

\subsubsection{Motivation}

Every SCOT experiment so far used SA-GAT as the backbone. The hypothesis was that maybe SA-GAT's skeleton-aware modifications don't play nicely with SCOT's learned prototypes --- type-pair attention, skeleton-relative position encoding, and same-type repulsion heads all make the embeddings more type-aware, which could conflict with SCOT's typed slot structure in subtle ways. A standard GATv2 with no skeleton modifications might give SCOT more flexibility.

The SCOT-KI experiment already suggested this wasn't the main issue (the assignment mechanism is the bottleneck, independent of embeddings) but we had never actually fine-tuned standard GAT + SCOT on COCO data. This section runs that experiment to close the hypothesis.

\subsubsection{Setup}

Loaded the pretrained standard GAT + SCOT checkpoint (synthetic-only, 0.923 synthetic / 0.862 COCO). Fine-tuned on COCO train2017 for 20 epochs at LR $1 \times 10^{-4}$ with the same infrastructure used for the SA-GAT fine-tune. SCOT head trained jointly (not frozen) so we can directly compare against the SA-GAT joint fine-tune.

\subsubsection{Results}

Best validation epoch was epoch 8 (val PGA 0.916). Test results:

\begin{tabular}{lccc}
\hline
\textbf{Method} & \textbf{Synthetic} & \textbf{COCO (GT kps)} & \textbf{End-to-end} \\
\hline
kNN (standard GAT, SCOT-trained) & 0.895 & 0.926 & 0.916 \\
\textbf{SCOT head (standard GAT)} & \textbf{0.939} & \textbf{0.881} & --- \\
COP-Kmeans (on these embeddings) & --- & --- & 0.934 \\
THA (on these embeddings) & --- & --- & 0.938 \\
\hline
\end{tabular}

\vspace{0.5em}

Comparison against the SA-GAT version of the same experiment:

\begin{tabular}{lcc}
\hline
\textbf{Variant} & \textbf{COCO SCOT} & \textbf{COCO kNN} \\
\hline
Standard GAT + SCOT (synthetic only) & 0.862 & 0.879 \\
Standard GAT + SCOT (COCO fine-tune, joint) & 0.881 & 0.926 \\
SA-GAT + SCOT (synthetic only) & 0.863 & 0.882 \\
SA-GAT + SCOT (COCO fine-tune, joint) & 0.884 & 0.946 \\
SA-GAT (no SCOT) + COP-Kmeans (COCO fine-tune) & --- & 0.957 \\
\hline
\end{tabular}

\subsubsection{Analysis}

\textbf{1. The hypothesis is closed: backbone choice does not matter for SCOT.} Standard GAT + SCOT lands at 0.881 COCO. SA-GAT + SCOT lands at 0.884 COCO. The 0.3-point difference is within noise. SCOT performance is determined by SCOT itself, not by which backbone feeds it. The skeleton-aware modifications in SA-GAT do not conflict with the SCOT loss --- they just don't help SCOT either.

\textbf{2. SA-GAT does produce better embeddings, but SCOT cannot capitalise on them.} On the joint COCO fine-tune, kNN reaches 0.946 with SA-GAT vs 0.926 with standard GAT --- SA-GAT adds 2 points of embedding quality. But SCOT scores 0.884 vs 0.881 on those same embeddings: the embedding improvement does not propagate through the SCOT head. This is consistent with SCOT-KI's behaviour on different embedding sources (HRNet features and SA-GAT both end up below COP-Kmeans/THA by similar margins).

\textbf{3. The SCOT loss costs $\sim$3 points of embedding quality regardless of backbone.} On standard GAT, kNN drops from the SA-GAT-only-fine-tune ceiling (0.957) to 0.926 when the SCOT loss is added --- a 3.1-point cost. On SA-GAT, the same comparison gives 0.957 $\to$ 0.946, a 1.1-point cost. Either way, training with SCOT actively degrades the embeddings the kNN/COP-Kmeans/THA methods would otherwise produce. This is the corruption effect we identified in the original SCOT joint-training experiment, now confirmed on a different backbone.

\textbf{4. Downstream methods (COP-Kmeans, THA) inherit the embedding degradation.} On the SCOT-trained standard GAT embeddings, end-to-end COP-Kmeans is 0.934 and THA is 0.938 --- both significantly below the SA-GAT-only fine-tune end-to-end results (0.955 each). The SCOT training pulled the embeddings into a worse-for-grouping region of feature space, and every downstream method pays the cost.

\textbf{5. The SCOT diagnosis is now bulletproof.} We have tested SCOT under:
\begin{itemize}
    \item 7 architectural variants (vanilla, residual, adaptive, unbalanced, dustbin, k-means init)
    \item 3 training strategies (joint synthetic, joint COCO fine-tune, frozen GAT)
    \item 2 backbone architectures (SA-GAT, standard GAT)
    \item 2 embedding sources (our contrastive embeddings, HigherHRNet's backbone features)
\end{itemize}
Across every combination, SCOT lands in the 0.83--0.88 COCO range while COP-Kmeans/THA reach 0.91--0.97. The architectural limitations of SCOT --- fixed prototypes, balanced Sinkhorn transport, soft-to-hard argmax conversion --- are the consistent bottleneck. No training strategy or embedding source rescues it.


\subsection{SCOT on HigherHRNet Backbone Features --- Diagnostic Experiment}

\subsubsection{Motivation}

Our SCOT analysis concluded that the architecture is the bottleneck: fixed prototypes don't adapt, balanced Sinkhorn transport assumes uniform joint visibility, and the ranking SCOT-KI $<$ THA $\approx$ COP-Kmeans is consistent across training conditions. But there was one alternative explanation we had not ruled out: what if the SA-GAT embeddings themselves are the problem? Perhaps a fundamentally different embedding source --- one not trained with contrastive loss, not derived from our GNN pipeline --- would unlock SCOT's potential.

The ideal test is to feed SCOT embeddings from a model that:
\begin{enumerate}
    \item Was trained on massive real COCO data (not 400 synthetic images)
    \item Produces features explicitly useful for grouping
    \item Is completely independent of our SA-GAT pipeline
\end{enumerate}

HigherHRNet~\cite{cheng2020higherhrnet} fits all three criteria. Its 32-dimensional backbone features are the raw features that both the heatmap head \emph{and} the AE tag grouping head read from --- trained end-to-end on 118k COCO images for detection and grouping. If any non-SA-GAT embedding should work with SCOT, it's these.

Note: HigherHRNet's AE tags are 1-dimensional scalars per keypoint, which cannot be used directly as SCOT embeddings (distance-based cost matrices need multi-dimensional vectors). The 32-dim backbone features are the closest useful alternative.

\subsubsection{Setup}

Registered a PyTorch forward hook on HigherHRNet's \texttt{final\_layers[0]} (the $\text{Conv2d}(32, 34, 1{\times}1)$ that outputs the heatmap + tag channels) to capture its 32-dimensional input feature map during inference. For each COCO val image, ran HigherHRNet detection, captured the backbone features, and bilinearly sampled the 32-dim feature vector at each detected keypoint location. Features were L2-normalised before use. Evaluated five grouping methods on these HigherHRNet backbone features: HigherHRNet's own AE grouping (baseline, uses the full grouping pipeline), kNN, COP-Kmeans, THA, and SCOT-KI.

\subsubsection{Results (2,165 COCO val images)}

\begin{tabular}{lcc}
\hline
\textbf{Method} & \textbf{HRNet features} & \textbf{SA-GAT (COCO ft)} \\
\hline
HigherHRNet AE grouping & 0.985 & --- \\
kNN & 0.789 & 0.957 \\
COP-Kmeans & 0.891 & 0.968 \\
THA & 0.884 & 0.969 \\
SCOT-KI & 0.825 & 0.926 \\
\hline
\end{tabular}

\subsubsection{Analysis}

\textbf{1. SA-GAT embeddings are strictly better than HigherHRNet backbone features for generic grouping.} kNN on SA-GAT gets 0.957 vs 0.789 on HRNet features --- a 17-point gap. Even though HRNet was trained on 118k real COCO images (300$\times$ more than our 400 synthetic images pre-training), its backbone features are not optimised for identity separation. They are optimised for detecting \emph{where} a joint is, not \emph{whose} joint it is. Our contrastive training explicitly optimises for identity, so the resulting embeddings are far more structured for grouping purposes. This is a positive result for the SA-GAT contribution.

\textbf{2. HigherHRNet's 0.985 grouping performance comes from the AE tag head, not the backbone features.} The tiny 561-parameter tag head is what makes HigherHRNet work. The backbone features alone barely separate identities. This contextualises the ``HigherHRNet is a full pipeline'' claim: 28.6M parameters do detection, 561 parameters do grouping, and the grouping only works because the tag head is trained end-to-end with a specific grouping loss. Swap the grouping method and the backbone features underperform.

\textbf{3. The SCOT-KI $<$ THA $\approx$ COP-Kmeans ordering holds on HRNet features.} On SA-GAT embeddings, the gap from SCOT-KI to COP-Kmeans is 0.926 to 0.968 (4.2 points). On HRNet features, it is 0.825 to 0.891 (6.6 points). The ranking is consistent and the gap is if anything larger with different embeddings. This is \textbf{definitive diagnostic evidence that SCOT's Sinkhorn assignment mechanism is the architectural bottleneck} --- independent of embedding source. The balanced transport assumption and argmax-from-soft-transport degradation happen regardless of what feature space the cost matrix is computed in.

\textbf{4. This closes off the ``embedding source'' hypothesis.} We have now tested SCOT-KI with: (a) synthetic-trained SA-GAT, (b) COCO-fine-tuned SA-GAT, and (c) HigherHRNet backbone features. All three produce SCOT-KI below COP-Kmeans by a consistent margin. No embedding source rescues SCOT. The paper 2 narrative is now rock-solid: ``we systematically tested SCOT across embedding sources and training strategies; the architectural limitations of fixed prototypes and balanced Sinkhorn transport are the reason it underperforms constraint-based assignment.''

\textbf{5. Side finding: COP-Kmeans and THA are robust to embedding quality.} On SA-GAT embeddings (kNN 0.957 $\to$ COP-Kmeans 0.968) the type constraint adds 1.1 points. On the much weaker HRNet features (kNN 0.789 $\to$ COP-Kmeans 0.891), the constraint adds 10.2 points. The worse the embeddings, the more the structural prior contributes. This suggests that type constraints are a particularly useful inductive bias when embedding quality is lower --- useful context for practitioners who can't fine-tune their detector.


\subsection{Visual Feature Augmentation Feasibility (Towards a Paper 3)}

\subsubsection{Motivation}

The MEng and PEng contributions so far are a \emph{grouping module} --- SA-GAT operates on already-detected keypoints and produces person assignments. It is modular: it can plug into any keypoint detector. But it is not an end-to-end pose estimation system. HigherHRNet, DEKR, ED-Pose, and similar methods all couple detection and grouping in a single trainable pipeline.

A potential third paper would close this gap by showing that SA-GAT can act as the grouping component in a full pose estimation system, where a lightweight backbone produces both (a) keypoint detections \emph{and} (b) per-keypoint visual features that SA-GAT consumes alongside positions. The contribution would shift from ``a grouping algorithm'' to ``a complete modular pose estimation pipeline with a learned grouping module that uses visual context.''

The first question to answer before committing to this direction is whether visual features actually help grouping. SA-GAT currently sees only positions and joint types --- no image context, no appearance information. If visual features add little, the contribution is marginal and not worth the engineering investment.

\subsubsection{Setup}

I built a feasibility test (\texttt{eval\_mobilenet\_features.py}) that:
\begin{enumerate}
    \item Detects keypoints with HigherHRNet (existing pipeline)
    \item Runs MobileNetV2 (ImageNet-pretrained) on the same image and captures an intermediate feature map
    \item Bilinearly samples the feature vector at each detected keypoint location
    \item Compares two embedding versions on the same scene:
    \begin{itemize}
        \item \textbf{pos}: $[x_{\text{norm}}, y_{\text{norm}}, \text{type one-hot}]$ (19-dim baseline)
        \item \textbf{pos+feat}: above concatenated with the MobileNet feature vector
    \end{itemize}
    \item Runs kNN, COP-Kmeans, and THA on both versions
\end{enumerate}

The test uses raw concatenated features without any learned processing --- it isolates the question ``do visual features carry useful identity signal?'' from ``can SA-GAT learn to use them?'' If the raw signal is useful, training SA-GAT to exploit it should amplify the gain.

I tested two MobileNetV2 layers:
\begin{itemize}
    \item \textbf{Mid} (layer 13, 96-dim, $/16$ resolution): mid-level features (textures, parts)
    \item \textbf{Deep} (layer 17, 320-dim, $/32$ resolution): high-level semantic features (object parts, appearance)
\end{itemize}

MobileNetV2 was chosen for being lightweight ($\sim$3.5M params), pretrained on ImageNet via torchvision (no Google Drive hunting), and general-purpose (not specialised for keypoint detection, so its features encode visual identity rather than localisation).

\subsubsection{Results (480 COCO val images)}

\textbf{Deep layer (320-dim, semantic):}

\begin{tabular}{lccr}
\hline
\textbf{Method} & \textbf{Position only} & \textbf{+ MobileNet} & \textbf{Diff} \\
\hline
kNN & 0.691 & 0.758 & $+0.067$ \\
COP-Kmeans & 0.857 & 0.886 & $+0.029$ \\
THA & 0.871 & 0.877 & $+0.006$ \\
\hline
\end{tabular}

\vspace{0.5em}

\textbf{Mid layer (96-dim, less semantic):}

\begin{tabular}{lccr}
\hline
\textbf{Method} & \textbf{Position only} & \textbf{+ MobileNet} & \textbf{Diff} \\
\hline
kNN & 0.691 & 0.705 & $+0.014$ \\
COP-Kmeans & 0.857 & 0.817 & $-0.040$ \\
THA & 0.871 & 0.797 & $-0.073$ \\
\hline
\end{tabular}

\subsubsection{Analysis}

\textbf{1. Deep semantic features carry identity signal.} The +6.7-point gain on kNN with deep features is significant --- visual context is genuinely informative for telling people apart. People look like themselves, and MobileNet's high-level features capture appearance differences that pure positions cannot.

\textbf{2. Mid-level features hurt rather than help.} Less semantic features (textures, edges) add noise instead of signal. They encode local visual properties that are similar between people in the same scene (lighting, background) rather than identity-distinguishing features. The deeper the feature, the more useful for grouping.

\textbf{3. The improvement diminishes as the grouping method gets stronger.} kNN gains 6.7 points, COP-Kmeans gains 2.9, THA gains only 0.6. The structural priors (type constraints, Hungarian matching) already do most of the work --- adding visual features mostly helps the weaker baselines catch up. By the time we have THA, the marginal value of visual features is barely above noise.

\textbf{4. Raw concatenation is unfair to the visual features.} This test uses raw concatenated features without any learned processing. SA-GAT could in principle learn to weight visual features differently per joint type, attend to neighbours' features, and combine them with positions in non-trivial ways. The test answers ``is the signal there?'' (yes for kNN, marginal for THA) but not ``can SA-GAT exploit it?'' which requires actual training.

\textbf{5. Feasibility verdict: marginal but worth a deeper test.} The +0.6 point gain on THA is too small to justify a full Paper 3 effort by itself. However, the +6.7 gain on kNN proves the signal exists. The question now is whether a trained SA-GAT v2 (with feature input) can extract more value from the visual features than raw concatenation does. A small training experiment on cached MobileNet features could answer this in 4--6 hours.

\subsubsection{Decision and next step}

Before committing to Paper 3, the next step is to train an SA-GAT v2 that accepts visual features as additional node input, on cached MobileNet (or HRNet) features from COCO train2017, for a few epochs. If the trained version pushes end-to-end PGA past 0.955 (the current SA-GAT + THA ceiling), Paper 3 is real. If it lands in the same 0.95--0.96 range, the visual features are not worth the architectural complexity and we focus on Papers 1 and 2.

This is the cheapest way to get a definitive answer before committing weeks of work to a full end-to-end pipeline.


\subsection{SA-GAT v2 --- Training with Visual Features (Actual Results)}

\subsubsection{Implementation}

I built the full pipeline:
\begin{itemize}
    \item \texttt{cache\_mobilenet\_features.py} --- runs MobileNetV2 (ImageNet-pretrained, layer 17, 320-dim) on every COCO train2017 and val2017 image, samples features at GT keypoint locations, saves a \texttt{.pt} file per image. 54,564 train / 2,258 val samples cached in $\sim$6 minutes on GPU.
    \item \texttt{cached\_features\_dataset.py} --- loads the cached \texttt{.pt} files and yields ready-to-use PyG \texttt{Data} objects with a \texttt{features} attribute on each node, bypassing the standard adapter/preprocessor pipeline entirely.
    \item \texttt{sa\_gat\_v2.py} --- a new SA-GAT variant that accepts a per-node visual feature vector. The features are projected through a $320 \to 32$ linear + LayerNorm + GELU block and concatenated with the standard positional + type embedding before the first GAT layer. All three skeleton-aware modifications (type-pair attention, position encoding, repulsion heads) are unchanged.
    \item Trainer and evaluator updates to support the SA-GAT v2 code path through optional config fields (fully backwards-compatible with existing configs).
\end{itemize}

Total SA-GAT v2 parameters: \textbf{197,280} --- only $\sim$27K more than the original SA-GAT (170,560), almost all of that in the $320 \to 32$ visual projection block.

\subsubsection{Training}

Trained from scratch on cached COCO train2017 features (54,564 samples) for 20 epochs at LR $1 \times 10^{-3}$, cosine decay to $1 \times 10^{-5}$, contrastive loss only (no grouping head). Validation on cached COCO val2017 features.

\begin{center}
\begin{tabular}{lc}
\hline
\textbf{Epoch} & \textbf{Val kNN PGA (COCO)} \\
\hline
2 & 0.960 \\
4 & 0.965 \\
8 & 0.969 \\
12 & 0.972 \\
16 & 0.972 \\
20 & 0.973 \\
\hline
\end{tabular}
\end{center}

The loss kept dropping through epoch 20 (0.081 $\to$ 0.034) and val PGA was still rising, so I trained for another 30 epochs (50 total) starting from the epoch-20 checkpoint at a halved learning rate ($5 \times 10^{-4}$).

\paragraph{Extended training (50 epochs total).}

\begin{center}
\begin{tabular}{lc}
\hline
\textbf{Epoch} & \textbf{Val kNN PGA (COCO)} \\
\hline
20 (end of initial run) & 0.9732 \\
30 & 0.9735 \\
40 & 0.9751 \\
50 & \textbf{0.9756} \\
\hline
\end{tabular}
\end{center}

The extended run added only $+0.002$ validation PGA over 30 more epochs. Val PGA plateaued around epoch 22 with tiny fluctuations afterwards --- the model is effectively converged at this scale. Continued training is not going to close the gap to HigherHRNet.

\subsubsection{Results (50-epoch final run)}

All results below are from the best checkpoint of the 50-epoch extended run.

\paragraph{COCO val with GT keypoints.}

\begin{tabular}{lcc}
\hline
\textbf{Method} & \textbf{SA-GAT (no features)} & \textbf{SA-GAT v2 (with features)} \\
\hline
kNN & 0.957 & \textbf{0.968} \\
COP-Kmeans & 0.968 & (not tested) \\
THA & 0.969 & (not tested) \\
\hline
\end{tabular}

\vspace{0.5em}

\paragraph{End-to-end with HigherHRNet detection (2,166 COCO val images).}

\begin{tabular}{lccc}
\hline
\textbf{Method} & \textbf{v1 (no features)} & \textbf{v2 (50 epochs)} & \textbf{$\Delta$} \\
\hline
SA-GAT + kNN & 0.942 & 0.958 & $+0.016$ \\
SA-GAT + COP-Kmeans & 0.955 & 0.963 & $+0.008$ \\
SA-GAT + THA & 0.955 & \textbf{0.964} & $+0.009$ \\
\hline
HigherHRNet AE & 0.985 & 0.985 & --- \\
\hline
\end{tabular}

\vspace{0.5em}

\paragraph{20 vs 50 epoch comparison.} The extended run gained essentially nothing end-to-end (THA: $0.9631 \to 0.9637$, a $+0.0006$ improvement over 30 additional epochs). This confirms the model is at its ceiling --- more training will not close the gap to HigherHRNet's $0.985$.

\vspace{0.5em}

Visual features helped every method, most for kNN ($+1.6$) and less for the structurally-aware methods ($+0.8$--$+0.9$ for COP-Kmeans and THA). The signal is real but small by the time the structural priors have done their work.

\paragraph{Synthetic test set.}

SA-GAT v2 dropped to 0.909 synthetic PGA (kNN). This is expected: the model was trained exclusively on cached COCO features and never saw synthetic data, and MobileNet's ImageNet-pretrained features on synthetic imagery are out-of-distribution. The synthetic numbers are not meaningful for this model.

\subsubsection{Honest analysis --- this is not a clear victory}

\textbf{1. We still do not beat HigherHRNet.} The end-to-end PGA with SA-GAT v2 + THA is 0.964 (after 50 epochs), compared to HigherHRNet's own associative embedding grouping at 0.985. We are \textbf{2.1 points behind} on the same detected keypoints. Training 30 more epochs on top of the initial 20 added essentially nothing ($+0.0006$), confirming we have hit the ceiling of what SA-GAT v2 can extract from this architecture and data.

\textbf{2. And we added more compute to do worse.} The total model size increased:
\begin{itemize}
    \item HigherHRNet alone: 28.6M parameters, 0.985 end-to-end PGA
    \item HigherHRNet + MobileNetV2 + SA-GAT v2: 30.8M parameters, 0.964 end-to-end PGA
\end{itemize}
We added 2.2M parameters (MobileNetV2 backbone) and 197K more (SA-GAT v2), a 7.7\% increase in total parameters, to achieve a 2.1-point drop in accuracy. This is the opposite of progress by any conventional metric.

\textbf{3. The uncomfortable conclusion: HigherHRNet's AE head is a better grouping module than SA-GAT v2 for this problem.} HigherHRNet's AE tag head is 561 parameters, trained end-to-end with the backbone on 118k real images, and achieves 0.985. SA-GAT v2 is 197K parameters, trained on the same backbone's detections plus additional visual features for 50 epochs, and achieves 0.964. The simpler, jointly-trained head wins.

\textbf{4. What the experiment \emph{does} show positively.}
\begin{itemize}
    \item \textbf{Visual features do help SA-GAT.} The v2 $\to$ v1 improvement (+0.8 end-to-end) is a real finding: adding appearance context to positional GNN reasoning measurably improves grouping. This validates the architectural intuition.
    \item \textbf{The modular approach works in principle.} SA-GAT v2 trained purely on cached features, with no joint training with the detector, reaches 0.963 --- within 2.2 points of a jointly-trained system. That's a small number if the story becomes ``detector-agnostic modular grouping'' rather than ``we beat HigherHRNet.''
    \item \textbf{Feature caching is a practical efficiency win.} 90-minute training on cached features, no detector retraining, portable to new detectors. If the goal is system modularity rather than accuracy maxxing, this is a valid contribution.
\end{itemize}

\textbf{5. What the experiment does \emph{not} show.} We have not demonstrated:
\begin{itemize}
    \item That SA-GAT v2 works with a different detector (we only tested with HigherHRNet)
    \item That it outperforms HigherHRNet's AE head in any scenario
    \item That visual features help more with weaker detectors (would be the obvious follow-up)
\end{itemize}

\textbf{6. Honest framing for Paper 3 (if it happens).} The paper cannot claim ``better than HigherHRNet.'' It would need to claim something like: ``a modular grouping component that achieves 93\% of an end-to-end system's accuracy while being detector-agnostic.'' That's a less exciting story, and whether it's publishable as a stand-alone paper is debatable. More likely, these results belong in the dissertation as an experimental finding (``we investigated end-to-end training with cached visual features and found that the jointly-trained AE baseline remains stronger for a single fixed detector'') rather than as a third paper.

\subsubsection{What to do next}

Three realistic options:

\textbf{Option A: Drop the Paper 3 idea.} Keep the two papers planned (SA-GAT + THA paper, SCOT investigation paper). The SA-GAT v2 work goes into the dissertation as an honest investigation with a negative framing. This is probably the right call.

\textbf{Option B: Pivot Paper 3 to a detector-agnostic story.} Train SA-GAT v2 to work with a second backbone (e.g.\ MMPose's ViTPose-tiny or a ResNet keypoint detector) and demonstrate that the same SA-GAT v2 works across both. The contribution becomes ``one grouping module, multiple backbones'' --- a weaker but still real story. Realistic timeline: 2--3 weeks of integration work.

\textbf{Option C: Find a setting where the AE head fails.} Test HigherHRNet's AE on genuinely difficult scenarios (high occlusion, large K, unusual poses) and see if SA-GAT v2 holds up better. If so, the paper claim becomes ``robust to occlusion'' rather than ``more accurate.'' Requires identifying a real weakness, which isn't guaranteed.

The honest recommendation is Option A. SA-GAT v2 is an interesting experiment that confirms visual features help our positional GNN, but it does not beat the baseline it was designed to compete with. That is a valid finding for the dissertation but not a strong standalone paper.


\subsection{Full Project Summary --- MEng and PEng Contributions}

This section provides a complete overview of all methods investigated, results achieved, and how the work splits across the two degrees. Written for the professor to assess scope and sufficiency.

\subsubsection{The conceptual split between the two degrees}

The two degrees address two different framings of the same underlying problem.

\paragraph{MEng --- grouping as a clustering problem, so make the embeddings as good as possible.}

The MEng starts from the standard view of bottom-up pose grouping: produce an embedding per keypoint, then cluster them into people. Everything in the MEng is in service of making those embeddings as good as possible for clustering. I started with a vanilla GAT and tried learned grouping heads on top (slot attention, graph partitioning, DEC) --- all of them failed to beat plain k-means on the same embeddings. I then built SA-DMoN and ran ten experiments attempting modularity-based clustering --- all of them plateaued around 0.77--0.81 because the modularity objective is fundamentally misaligned with pose grouping (the kNN graph encodes spatial proximity, but communities are defined by identity).

The lesson from those failures was that \emph{the embeddings are the lever}, not the grouping head on top. SA-GAT --- with three skeleton-aware modifications (type-pair attention bias, skeleton-relative position encoding, same-type repulsion heads) --- is the embedding network that comes out of that lesson. With well-structured contrastive embeddings, even plain kNN (0.957 COCO end-to-end after fine-tuning) matches or exceeds every learned grouping head tested.

The MEng thesis is: \textbf{multi-person pose grouping reduces to an embedding problem; SA-GAT produces embeddings good enough that simple clustering outperforms every learned grouping head I tried.}

\paragraph{PEng --- grouping is not a clustering problem, it is an assignment problem.}

The PEng takes SA-GAT's embeddings as a given and pivots conceptually. The starting observation is COP-Kmeans: constrained k-means with a must-not-link constraint that no two left elbows can belong to the same person. COP-Kmeans beats every learned grouping head the MEng tested using zero learned parameters --- but the reason it works is not that it is ``a better cluster algorithm''. The must-not-link constraint is not a clustering property; it is an assignment property. Each person has exactly 17 typed slots (one per COCO joint type), and the grouping task is to assign detected keypoints to slots under the hard constraint that each slot receives at most one keypoint. Viewed this way, pose grouping is structurally closer to the linear assignment problem or optimal transport than it is to k-means clustering.

This reframing is where the rest of the PEng comes from. SCOT formalises the assignment view: the same hard type constraint that COP-Kmeans enforces greedily, embedded in a principled differentiable optimal transport formulation with typed skeleton slots ($K \times 17$), infinite cost masking for type violations, and Sinkhorn iterations.

SCOT underperforms. Through systematic ablation (seven variants, three training strategies, two backbones, two embedding sources) I diagnosed two architectural weaknesses: fixed learned prototypes that do not adapt to scene diversity, and balanced Sinkhorn transport that assumes uniform joint visibility. SCOT-KI (k-means initialised prototypes) fixes the first problem and recovers 6.5 points. THA (Type-wise Hungarian Assignment) fixes both: scene-adaptive k-means centroids plus per-joint-type optimal bipartite matching via Hungarian, which eliminates the balanced transport assumption. THA beats COP-Kmeans on every condition tested.

The PEng thesis is: \textbf{multi-person pose grouping is structurally an assignment problem with skeleton constraints, not a clustering problem; the right algorithm is type-decomposed optimal bipartite matching.}

\paragraph{The relationship between the two.}

The MEng builds the embedding machinery under a clustering interpretation and finds that constraint-based clustering (COP-Kmeans) beats learned heads. The PEng takes SA-GAT as a given, reframes the grouping task as structured assignment rather than clustering, and investigates assignment-based methods (SCOT, SCOT-KI, THA) that directly target the type constraints as a structural property of the problem. The two degrees share the embedding network and the dataset infrastructure but diverge in how they frame what ``grouping'' actually is.


\subsubsection{Infrastructure contributions (shared)}

\begin{itemize}
    \item \textbf{Unified data pipeline.} Adapter pattern (VirtualAdapter, CocoAdapter) producing identical output formats. Single preprocessor builds PyG kNN graphs from either source. Depth toggling via config (MiDaS~\cite{ranftl2020midas} skipped when disabled).
    \item \textbf{Synthetic data generator.} Blender-based procedural scene generation producing 400 training / 100 validation / 100 test scenes with perfect GT keypoints and person labels.
    \item \textbf{COCO integration.} Full COCO train2017/val2017 support with GT keypoint loading, per-joint depth estimation (optional), and fine-tuning support in the trainer.
    \item \textbf{Training infrastructure.} Pydantic YAML configs, GUID-based output directories (no checkpoint overwrites), latest symlinks, config copying for reproducibility.
    \item \textbf{Evaluation infrastructure.} PGA metric (Hungarian-matched joint assignment accuracy), per-joint breakdowns, end-to-end pipeline with HigherHRNet-W32-512~\cite{cheng2020higherhrnet} detection, COCO/synthetic comparison.
    \item \textbf{K estimation head.} Small MLP trained on frozen embeddings to predict person count. 79.6\% exact / 96\% off-by-one on COCO val after COCO training.
\end{itemize}


\subsubsection{MEng: embedding network (grouping-as-clustering)}

The MEng work lives under the clustering framing. It covers (a) the failed learned grouping heads that motivated the embedding-focused approach, (b) SA-DMoN (the most time-intensive negative result), and (c) SA-GAT itself as the embedding network that comes out of those failures.

\paragraph{Embedding network investigations.}

\begin{tabular}{llcc}
\hline
\textbf{Method} & \textbf{Description} & \textbf{Synth PGA} & \textbf{COCO PGA} \\
\hline
Standard GAT~\cite{brody2022gatv2} & 2-layer GATv2, contrastive loss & 0.901 & 0.879 \\
Standard GAT (with depth) & Same + MiDaS depth & 0.994 & 0.838 \\
SA-GAT (type-pair only) & + category attention bias & 0.909 & --- \\
SA-GAT (pos enc only) & + skeleton-relative position & 0.908 & --- \\
SA-GAT (repulsion only) & + same-type repulsion heads & 0.897 & --- \\
SA-GAT (full) & All three modifications & 0.910 & 0.882 \\
SA-GAT (COCO fine-tuned) & Fine-tuned 20 epochs on COCO & 0.928 & 0.957 \\
\hline
\end{tabular}

\vspace{0.5em}

Key finding: removing depth improved COCO transfer by +4.1 points (MiDaS domain gap). SA-GAT's three modifications each contribute independently. COCO fine-tuning was the single largest improvement (+7.5 points kNN on COCO).

\paragraph{Learned grouping head attempts (all failed to beat plain kNN).}

All tested on standard GAT embeddings with GT K. Slot attention and graph partitioning were trained with depth (these predate the no-depth standard).

\begin{tabular}{llcc}
\hline
\textbf{Method} & \textbf{Type} & \textbf{Synth} & \textbf{COCO} \\
\hline
kNN (k-means baseline) & No head & 0.901 & 0.879 \\
Slot attention (depth) & Learned head & 0.858 & 0.788 \\
Graph partitioning (depth) & Edge classifier & 0.954 & 0.752 \\
DEC & Deep clustering & 0.654 & --- \\
\hline
\end{tabular}

\vspace{0.5em}

Every learned head failed to outperform plain kNN clustering on the same embeddings. The contrastive embeddings were already good enough that learned assignment added nothing. This was the first evidence that the embeddings, not the grouping head, are the real lever.

\paragraph{SA-DMoN: modularity-based clustering (10 experiments).}

\begin{tabular}{llc}
\hline
\textbf{Experiment} & \textbf{Variant} & \textbf{Synth PGA} \\
\hline
Exp 3 & Vanilla DMoN~\cite{tsitsulin2023dmon}, no depth & 0.785 \\
Exp 4 & SA-DMoN v1, learnable sigma & 0.807 \\
Exp 5 & SA-DMoN v1, sigma clamped [0.05, 0.5] & 0.812 \\
Exp 6 & SA-DMoN v1, sigma clamped, $\lambda_{\text{spectral}}=1.0$ & 0.799 \\
Exp 7 & SA-DMoN v1, sigma = median distance & 0.770 \\
Exp 8 & SA-DMoN v1, spectral loss removed & 0.806 \\
Exp 9 & SA-DMoN v2, feature adjacency $\lambda=1.0$ & 0.781 \\
Exp 10 & SA-DMoN v2, feature adjacency $\lambda=0.1$ & 0.797 \\
Exp 11 & SA-DMoN v2, entropy type loss only & 0.806 \\
Exp 12 & SA-DMoN v2, full (both fixes) & 0.801 \\
\hline
\end{tabular}

\vspace{0.5em}

All SA-DMoN variants converge to 0.77--0.81 regardless of formulation. Root cause: the modularity objective is fundamentally misaligned with pose grouping. The kNN graph encodes spatial proximity, but communities are defined by identity. The spatial null model ($R_{\text{pose}}$) either cancels with the spatial adjacency ($A$) or is too weak/strong. This was the most time-intensive investigation and the clearest negative result --- it exhaustively eliminated modularity-based clustering as a viable approach.

\paragraph{MEng conclusion.} With SA-GAT producing well-structured contrastive embeddings, grouping-as-clustering is effectively solved. Simple kNN already gets 0.957 on COCO after fine-tuning; any sensible clustering method on top reaches 0.95+. The MEng paper's claim is that pose grouping is an embedding problem and SA-GAT is the right embedding network for it.


\subsubsection{PEng: grouping-as-assignment (COP-Kmeans $\to$ SCOT $\to$ THA)}

The PEng work takes SA-GAT as a given and reframes the problem. Hard type constraints (no two left elbows per person) are not a clustering property --- they are an assignment property. This reframing starts with COP-Kmeans, formalises through SCOT, and ends with THA.

\paragraph{Constraint-based assignment variants (on SA-GAT embeddings).}

\begin{tabular}{llcc}
\hline
\textbf{Method} & \textbf{Mechanism} & \textbf{Synth} & \textbf{COCO} \\
\hline
COP-Kmeans~\cite{wagstaff2001copkmeans} & Greedy constrained k-means (must-not-link) & 0.913 & 0.910 \\
Agglomerative & Constrained bottom-up merging & 0.837 & 0.887 \\
Spectral & Zero same-type affinities & 0.915 & 0.894 \\
\hline
\end{tabular}

\vspace{0.5em}

COP-Kmeans is the first clear signal that the constraint is the contribution --- it takes the same embeddings the learned heads use and beats them with no learned parameters. The must-not-link constraint is doing the work. This is the pivot moment: type constraints are an \emph{assignment} construct, not a clustering construct. The remainder of the PEng investigates that reframing properly.

\paragraph{SCOT: skeleton-constrained optimal transport.}

\begin{tabular}{llcc}
\hline
\textbf{Variant} & \textbf{Training} & \textbf{Synth} & \textbf{COCO} \\
\hline
SCOT (fixed prototypes) & Joint, synthetic & 0.926 & 0.863 \\
SCOT (fixed prototypes) & Joint, COCO fine-tune & 0.948 & 0.884 \\
SCOT (fixed prototypes) & Frozen GAT, COCO & 0.838 & 0.861 \\
Residual SCOT & Joint, synthetic & 0.920 & 0.868 \\
Adaptive SCOT & Joint, synthetic & 0.919 & 0.854 \\
Unbalanced SCOT & Joint, synthetic & 0.894 & 0.836 \\
Dustbin SCOT & Joint, synthetic & 0.901 & 0.840 \\
SCOT-KI (k-means prototypes) & No training & --- & 0.926 \\
\hline
COP-Kmeans (for comparison) & No training & 0.938 & 0.968 \\
\hline
\end{tabular}

\vspace{0.5em}

SCOT's novel formulation (typed skeleton slots, infinite cost masking, Sinkhorn OT) achieves the highest synthetic PGA of any learned method (0.926). However, every SCOT variant underperforms COP-Kmeans on COCO by 4--13 points. Systematic diagnosis identified two causes: (1) fixed prototypes cannot represent the diversity of real scenes, and (2) balanced Sinkhorn transport assumes uniform joint visibility which doesn't hold with occlusion. SCOT-KI (scene-adaptive k-means prototypes) recovered 6.5 points but the remaining 4.2-point gap is in the Sinkhorn mechanism itself.

\paragraph{THA: type-wise Hungarian assignment (final method).}

Motivated by the SCOT analysis: replace both the fixed prototypes (with k-means centroids) and the Sinkhorn assignment (with per-type Hungarian matching).

\begin{tabular}{lcc}
\hline
\textbf{Method} & \textbf{Synth} & \textbf{COCO} \\
\hline
kNN & 0.928 & 0.957 \\
COP-Kmeans & 0.938 & 0.968 \\
\textbf{THA} & \textbf{0.965} & \textbf{0.969} \\
\hline
\end{tabular}

\vspace{0.5em}

THA beats COP-Kmeans on both datasets (+2.7 synthetic, +0.1 COCO) with zero learned parameters. End-to-end with HigherHRNet: 0.955 (matching COP-Kmeans). Novel formulation: 17 parallel bipartite matching problems connected through shared centroids with iterative refinement.

\paragraph{Other PEng investigations.}
\begin{itemize}
    \item \textbf{Skeleton-aware edge classifier} --- binary edge classification for same-person/different-person edges. Abandoned after GEC~\cite{maher2025gec} published the same approach in early 2025.
    \item \textbf{K estimation} --- eigenvalue gap (18.3\% exact), joint MLP training (corrupted embeddings), frozen backbone MLP (100\% synthetic, 79.6\% COCO).
    \item \textbf{Confidence-gated SCOT+kNN fusion} --- did not improve over SCOT alone.
    \item \textbf{SA-GAT v2 with MobileNet features} --- attempt at a third paper (end-to-end pipeline with visual features). Trained for 50 epochs, reached 0.964 end-to-end vs HigherHRNet's 0.985, so the paper was dropped. Visual features do help (+0.8 points) but not enough to beat a jointly-trained AE head. Documented as a negative result for the dissertation.
\end{itemize}

\paragraph{PEng conclusion.} The assignment framing pays off: COP-Kmeans establishes that the constraint is the contribution, SCOT formalises it as OT but underperforms due to fixed prototypes and balanced transport, and THA realises the framing in a form that beats every clustering baseline while using zero learned parameters. The PEng paper's claim is that pose grouping is structurally an assignment problem with skeleton constraints, and type-decomposed Hungarian matching is the right algorithm for it.


\subsubsection{Complete results table}

All methods, best numbers, on COCO val with GT keypoints and GT K (except where noted):

\begin{tabular}{llc}
\hline
\textbf{Method} & \textbf{Backbone} & \textbf{COCO PGA} \\
\hline
\multicolumn{3}{l}{\emph{Baselines:}} \\
kNN & Standard GAT & 0.879 \\
kNN & SA-GAT & 0.882 \\
kNN & SA-GAT (COCO ft) & 0.957 \\
\hline
\multicolumn{3}{l}{\emph{Learned grouping heads:}} \\
Slot attention & Standard GAT (depth) & 0.788 \\
Graph partitioning & Standard GAT (depth) & 0.752 \\
SA-DMoN (best, no depth) & Standard GAT & 0.817 \\
SCOT & Standard GAT & 0.862 \\
SCOT & SA-GAT & 0.863 \\
SCOT (COCO ft, joint) & SA-GAT & 0.884 \\
SCOT (COCO, frozen) & SA-GAT (COCO ft) & 0.861 \\
\hline
\multicolumn{3}{l}{\emph{Constraint-based grouping (no learned params):}} \\
COP-Kmeans & Standard GAT & 0.910 \\
COP-Kmeans & SA-GAT & 0.915 \\
COP-Kmeans & SA-GAT (COCO ft) & 0.968 \\
Agglomerative & Standard GAT & 0.887 \\
Spectral & Standard GAT & 0.894 \\
SCOT-KI & SA-GAT (COCO ft) & 0.926 \\
\textbf{THA} & \textbf{SA-GAT (COCO ft)} & \textbf{0.969} \\
\hline
\multicolumn{3}{l}{\emph{End-to-end (HigherHRNet detected keypoints):}} \\
HigherHRNet AE & HigherHRNet-W32 & 0.985 \\
SA-GAT + COP-Kmeans & SA-GAT (COCO ft) & 0.955 \\
SA-GAT + THA & SA-GAT (COCO ft) & 0.955 \\
\hline
\end{tabular}

\vspace{0.5em}

Total: 20+ distinct methods/configurations evaluated across synthetic and COCO datasets. 10 SA-DMoN experiments. 7 SCOT variants. 4 type-constrained clustering methods. SA-GAT ablation study (4 configurations). COCO fine-tuning. End-to-end pipeline with external detector. K estimation (3 approaches).


\subsection{Hyperbolic SA-GAT --- Lorentz-Model Embedding Space}

\subsubsection{Motivation}

A late-project experiment driven by a theoretical curiosity: skeletons are trees (kinematic chain from pelvis outward), and multi-person scenes are forests. Hyperbolic space has exponentially growing volume with radius, which exactly matches tree growth, whereas Euclidean space grows only polynomially. This makes hyperbolic geometry a provably better ambient space for tree-structured data at low distortion.

The hypothesis: if SA-GAT is currently spending attention capacity implicitly learning hierarchical structure, replacing the Euclidean output space with a hyperbolic one should encode that hierarchy for free. Two pre-registered hypotheses:

\begin{itemize}
    \item \textbf{H1:} hyperbolic beats Euclidean at matched embedding dimension.
    \item \textbf{H2:} hyperbolic matches Euclidean at significantly lower embedding dimension (e.g.\ 32 vs 128).
\end{itemize}

Realistic priors going in: COCO skeletons have only 17 joints and depth ~4, so the tree-distortion argument is asymptotic and probably weak at this scale. Expected outcomes in order of likelihood: no significant change, small positive gain, improvement at low dim only. Expected cost: $\sim$half a day if the numerics cooperate.

\subsubsection{Implementation}

Created \texttt{sa\_gat\_hyperbolic.py} alongside \texttt{sa\_gat.py} (non-invasive --- no changes to the existing SA-GAT code). Added \texttt{SAGATHyperbolicConfig} and registered it in \texttt{ExperimentConfig} as another optional field, following the same pattern as SA-GAT v2. The trainer, evaluator, and eval\_cop\_kmeans were all updated with a new branch that constructs the hyperbolic model and includes a hyperbolic-to-Euclidean mapping step for downstream grouping.

\textbf{Design: minimal-change hyperbolic embedding.} The internal SA-GAT attention layers remain Euclidean. Only the final output is projected onto the Lorentz hyperboloid via exponential map at the origin. This tests ``is a hyperbolic output metric better for pose grouping?'' without committing to a full hyperbolic attention rewrite.

The Lorentz model (a.k.a.\ hyperboloid model) was chosen over the Poincar\'e ball because it is more numerically stable, which several recent papers note. \texttt{geoopt}'s \texttt{geoopt.Lorentz(k=1.0)} provides the manifold with $\text{expmap}_0$, $\text{logmap}_0$, and geodesic distance.

\textbf{Contrastive loss in hyperbolic space.} The standard GATOnlyLoss uses cosine similarity on L2-normalised Euclidean embeddings. I wrote a new \texttt{HyperbolicContrastiveLoss} that computes pairwise geodesic distances on the Lorentz manifold and applies the standard margin formulation: pull same-person distances to zero, push different-person distances beyond a margin.

\subsubsection{Debugging issues encountered}

\textbf{Issue 1: NaN loss on the first forward pass.} Using \texttt{embeddings.unsqueeze(1).expand(n, n, -1)} to build pairwise tensors without \texttt{.contiguous()} broke geoopt's distance computation in two ways:
\begin{itemize}
    \item Diagonal distances (distance from a point to itself) came out as $\sim 5.5$ instead of $0$, indicating the broadcast was silently misaligning the Minkowski inner product calculation.
    \item Some off-diagonal pairs produced NaN because $\text{arccosh}$ received values slightly below 1 due to floating-point error on the misaligned tensors.
\end{itemize}
Fix: add \texttt{.contiguous()} calls on the expanded tensors and add a \texttt{torch.nan\_to\_num} clamp on the distance matrix for residual numerical edge cases. After the fix, diagonal distances were $\sim 10^{-4}$ (acceptable) and no NaN values appeared.

\textbf{Issue 2: Exploding distances at initialisation.} With the final LayerNorm producing tangent vectors of norm $\sim \sqrt{D} \approx 11$, the corresponding hyperboloid points were at geodesic distance $\sim 11$ from the origin, and pairwise distances were in the range $18$--$22$. The default contrastive margin of $1.0$ was meaningless at that scale (all negatives trivially satisfied the margin, no gradient). The initial loss was $\sim 20$ with pos and neg distances nearly identical around $20$.

Fix: added a learnable scalar \texttt{tangent\_scale} parameter that multiplies the tangent vector before the exponential map. Initialised to $1/\sqrt{D}$, which puts initial pairwise distances in the reasonable range $\sim 0.6$--$1.0$. The scale is a free parameter the optimiser can adjust during training. This is a standard trick in hyperbolic deep learning to keep the model away from the numerically dangerous regions of the manifold (far from origin, near the boundary of the Poincar\'e ball, etc.).

\textbf{Issue 3: Margin calibration.} The Euclidean SA-GAT uses a cosine margin of $0.5$, interpreted as ``negative pairs should have cosine similarity below $0.5$''. In hyperbolic space the units are different: the margin is a geodesic distance. After the scale fix, I set the margin to $2.0$ so negative pairs were pushed beyond $2\times$ the initial average distance. Training proceeded smoothly with this setting.

\textbf{Other minor fixes.} The \texttt{eval\_cop\_kmeans.py} and \texttt{evaluator.py} scripts originally only knew about standard GAT and SA-GAT; I extended both to handle the hyperbolic case. For grouping, the hyperbolic embeddings are mapped back to the tangent space at the origin via $\text{logmap}_0$, the extra time coordinate is dropped, and the result is L2-normalised so the existing Euclidean-based grouping methods (kNN, COP-Kmeans, THA) can operate on them directly. This is a principled choice: the tangent space at origin is the Euclidean linearisation of the manifold in the neighbourhood of origin, and distances in it approximate hyperbolic distances for small radii.

\subsubsection{Compromises made}

\begin{itemize}
    \item \textbf{Hybrid Euclidean/hyperbolic architecture.} The internal SA-GAT layers are Euclidean; only the final embedding is hyperbolic. A full hyperbolic attention mechanism (all layers on the manifold with M\"obius operations and Fr\'echet-mean aggregation) would be a more principled test but a much larger rewrite. If H1 had shown a large gain, the full rewrite would be the next step. With the modest gain observed, the hybrid is sufficient to answer the research question.
    \item \textbf{Grouping in tangent space.} Downstream clustering (kNN, COP-Kmeans, THA) operates on $\text{logmap}_0$ of the hyperboloid embeddings rather than on the hyperboloid itself. This approximation is good near the origin (where our \texttt{tangent\_scale} keeps most points) but distorts for far-from-origin points. A hyperbolic-metric variant of COP-Kmeans would be more faithful but also a much larger engineering effort.
    \item \textbf{No Riemannian optimiser.} Parameters are stored in Euclidean tangent space and optimised with standard AdamW. A Riemannian optimiser (RAdam/RSGD from geoopt) would be theoretically more correct when training parameters that live directly on the manifold, but our model has no such parameters --- the only manifold operation is expmap at the forward pass output --- so standard optimisers are fine.
\end{itemize}

\subsubsection{Results --- H1 (matched dim = 128, 50 epochs)}

\begin{tabular}{lcccc}
\hline
\textbf{Method} & \textbf{Synth kNN} & \textbf{Synth COP-KM} & \textbf{COCO kNN} & \textbf{COCO COP-KM} \\
\hline
SA-GAT Euclidean (baseline) & 0.910 & 0.925 & 0.882 & 0.915 \\
\textbf{SA-GAT Hyperbolic} & 0.907 & 0.916 & \textbf{0.895} & \textbf{0.924} \\
$\Delta$ & $-0.003$ & $-0.009$ & $\mathbf{+0.013}$ & $\mathbf{+0.009}$ \\
\hline
\end{tabular}

\vspace{0.5em}

\paragraph{Transfer gap (synthetic $\to$ COCO):}

\begin{tabular}{lcc}
\hline
\textbf{Variant} & \textbf{Transfer gap (kNN)} & \textbf{Transfer gap (COP-KM)} \\
\hline
Euclidean & $-2.8$ & $-1.0$ \\
\textbf{Hyperbolic} & $-1.2$ & $\mathbf{+0.8}$ \\
\hline
\end{tabular}

\vspace{0.5em}

The hyperbolic variant slightly underperforms Euclidean on in-distribution synthetic data ($-0.3$ to $-0.9$ points) but consistently outperforms on zero-shot COCO ($+0.9$ to $+1.3$ points). With COP-Kmeans, hyperbolic on COCO ($0.924$) is actually higher than on synthetic ($0.916$) --- a positive transfer gap, which is unusual for synthetic-trained models.

\subsubsection{Results --- H2 (lower dim = 32, 50 epochs)}

Same training setup as H1 but with \texttt{output\_dim: 32} (a $4\times$ reduction). All other hyperparameters held constant.

\begin{tabular}{lcccc}
\hline
\textbf{Method} & \textbf{Synth kNN} & \textbf{Synth COP-KM} & \textbf{COCO kNN} & \textbf{COCO COP-KM} \\
\hline
SA-GAT Euclidean (dim=128) & 0.910 & 0.925 & 0.882 & 0.915 \\
SA-GAT Hyperbolic (dim=128) & 0.907 & 0.916 & 0.895 & 0.924 \\
\textbf{SA-GAT Hyperbolic (dim=32)} & \textbf{0.918} & \textbf{0.930} & \textbf{0.899} & \textbf{0.926} \\
\hline
\end{tabular}

\vspace{0.5em}

\paragraph{The surprise.} H2 was expected to match Euclidean at best. Instead, \textbf{hyperbolic at dim=32 outperforms both hyperbolic at dim=128 and Euclidean at dim=128 on every metric}. Synthetic kNN: $0.918$ (H2) vs $0.910$ (Euclidean). COCO COP-Kmeans: $0.926$ (H2) vs $0.915$ (Euclidean).

Best validation PGA during training also climbed noticeably higher for H2 (0.882 vs 0.878 for H1). The val-PGA curve was still slowly improving at epoch 45 for H2 (last best-update), which is consistent with a tighter embedding space that the optimiser can exploit further.

\paragraph{Transfer gap:}

\begin{tabular}{lcc}
\hline
\textbf{Variant} & \textbf{Transfer gap (kNN)} & \textbf{Transfer gap (COP-KM)} \\
\hline
Euclidean (dim=128) & $-2.8$ & $-1.0$ \\
Hyperbolic (dim=128) & $-1.2$ & $+0.8$ \\
\textbf{Hyperbolic (dim=32)} & $-1.9$ & $-0.4$ \\
\hline
\end{tabular}

\vspace{0.5em}

The dim=32 model has a slightly larger transfer gap than dim=128 hyperbolic but still smaller than Euclidean, and it starts from a higher baseline on both datasets. So the transfer advantage is partially consumed by better in-distribution fit --- but the absolute COCO number is the highest of all three variants.

\subsubsection{Analysis}

\textbf{1. Both H1 and H2 succeeded --- H2 more strongly.} H1 (matched dim 128) showed a small transfer-advantage effect: $-0.3$ to $-0.9$ on synthetic, $+0.9$ to $+1.3$ on COCO. H2 (reduced dim 32) overtook \emph{both} Euclidean dim-128 and hyperbolic dim-128 on every single metric --- synthetic and COCO, kNN and COP-Kmeans. The theoretical argument (hyperbolic space needs less dimension to embed trees without distortion) appears to be empirically supported here, at small but consistent scale.

\textbf{2. Dim=32 hyperbolic is the best pre-COCO-fine-tune embedding tested.} SA-GAT Euclidean at dim=128 gets 0.915 COCO COP-Kmeans. SA-GAT Hyperbolic at dim=32 gets 0.926 on the same. That is $+1.1$ points, achieved with $1/4$ the embedding size. This is a concrete, defensible result: \textbf{same or better pose grouping accuracy at $4\times$ smaller embedding, by replacing Euclidean output with Lorentz-model hyperbolic output.}

\textbf{3. Why might dim=32 beat dim=128?} The theoretical claim is that hyperbolic space can embed trees at arbitrary low distortion in any dimension $\geq 2$, whereas Euclidean space needs $\Omega(\log n)$ dimensions for distortion bounded below a constant. At dim=128 in hyperbolic space, most of the ambient capacity is unused --- the optimiser has too much freedom and the regularisation implicit in LayerNorm $\to$ expmap does not prevent over-parameterisation. At dim=32 the capacity is closer to what the problem actually needs, and the geometry provides the structure the Euclidean model had to learn explicitly. This matches the standard finding in hyperbolic deep learning literature: lower-dimensional hyperbolic often outperforms higher-dimensional Euclidean.

\textbf{4. The positive transfer gap seen with H1 partially disappears for H2.} H1's most striking result was COCO COP-KM ($0.924$) beating synthetic ($0.916$). H2 gets the same COCO number ($0.926$) but a higher synthetic number ($0.930$), so the transfer gap returns to slightly negative ($-0.4$). H2 wins in absolute COCO accuracy but does not reproduce H1's ``better on OOD than ID'' result. The absolute numbers matter more for practical use, so H2 is still the preferred variant --- but H1's oddly-positive transfer gap is worth flagging as an interesting observation in itself.

\textbf{5. The compromises remain, and they still cap the strength of the conclusion.} Internal attention is still Euclidean; grouping is still done in tangent space. A full hyperbolic architecture might show larger effects (or might not --- the distortion argument is asymptotic and pose skeletons are tiny trees). We cannot claim ``hyperbolic is strictly better for pose grouping'' without a full rewrite; we can only claim ``moving the output to hyperbolic, with these specific choices, produces a consistent advantage.''

\textbf{6. This now supports a proper ablation row in the MEng paper.} Two concrete claims that survive scrutiny:
\begin{itemize}
    \item ``SA-GAT with Lorentz-model hyperbolic output embeddings at dim=32 matches or beats the Euclidean variant at dim=128 across all tested grouping methods and both datasets.''
    \item ``This is consistent with the theoretical prediction that hyperbolic space requires fewer dimensions to embed tree-structured data without distortion.''
\end{itemize}
Not a headline contribution but not trivial either. Worth 1--2 tables and half a page of discussion in the MEng paper as a principled efficiency result. Gets the thesis a sliver of theoretical grounding it did not previously have.

\subsubsection{What to do next}

Three realistic options:

\textbf{Option A: Stop here.} H2 gave a clear positive result that stands as an ablation row in the MEng paper. We have documented it properly. The marginal value of further hyperbolic experiments is low relative to other work that needs doing (paper writing, figure generation, thesis compilation).

\textbf{Option B: Train a COCO fine-tuned hyperbolic dim=32.} The Euclidean variant fine-tuned goes from $0.915 \to 0.968$ COCO COP-KM ($+5.3$ points). If hyperbolic dim=32 gets a similar jump, we would be at $\sim 0.980$ COCO --- which would be within 0.5 points of HigherHRNet's $0.985$. That is a genuinely thesis-worthy number. Cost: ~90 minutes of training + evaluation. Recommended.

\textbf{Option C: Full hyperbolic architecture.} Rewrite SA-GAT's attention to operate entirely on the manifold (M\"obius operations, Fr\'echet-mean aggregation). This is the principled version. Cost: a week or more. Probably not worth it unless an early reviewer explicitly asks for it.

The honest recommendation is Option B, then stop. It is the cheapest way to convert an ablation into a potentially headline result.

\subsubsection{Option B results --- COCO fine-tuned hyperbolic dim=32}

Ran the COCO fine-tune (Option B) with the same settings as the Euclidean fine-tune: load \texttt{outputs/train\_sa\_gat\_hyperbolic\_dim32/latest/best.pt}, fine-tune on COCO train2017 for 20 epochs at LR $1 \times 10^{-4}$, contrastive loss only, no grouping head.

The prediction ($\sim 0.980$ COCO, within 0.5 points of HigherHRNet) was wrong. The actual numbers came in lower. Full comparison across all four hyperbolic configurations plus the two Euclidean references:

\begin{tabular}{lcccc}
\hline
\textbf{Model} & \textbf{Synth kNN} & \textbf{COCO kNN} & \textbf{COCO COP-KM} & \textbf{COCO THA} \\
\hline
Euclidean dim=128 (synth only) & 0.910 & 0.882 & 0.915 & --- \\
Euclidean dim=128 (COCO ft) & 0.928 & \textbf{0.957} & \textbf{0.968} & \textbf{0.969} \\
Hyperbolic dim=128 (synth only) & 0.907 & 0.895 & 0.924 & --- \\
Hyperbolic dim=32 (synth only) & 0.918 & 0.899 & 0.926 & --- \\
Hyperbolic dim=32 (COCO ft) & 0.889 & 0.943 & 0.958 & 0.961 \\
\hline
\end{tabular}

\vspace{0.5em}

\paragraph{What the numbers say.}

\begin{itemize}
    \item \textbf{COCO fine-tuning works for the hyperbolic model.} Real gains on every COCO metric ($+4.4$ kNN, $+3.2$ COP-KM). Absolute numbers are strong in isolation (0.958 COP-KM, 0.961 THA).
    \item \textbf{Euclidean dim=128 COCO-ft still wins.} The hyperbolic variant trails by 0.8--1.4 points on every COCO metric after both are fine-tuned. The dim=32 advantage from the synthetic-only regime does not survive the fine-tune.
    \item \textbf{Hyperbolic synthetic performance degraded during fine-tuning.} Synth kNN dropped $0.918 \to 0.889$ ($-2.9$), while the Euclidean variant's synthetic went $0.910 \to 0.928$ (\emph{up} $+1.8$) during its own fine-tune. The hyperbolic model specialised more aggressively to COCO.
    \item \textbf{The transfer-advantage story partially reversed.} Post-fine-tune, the hyperbolic synth-to-COCO gap is $+5.4$ (COCO $>$ synth) while the Euclidean gap is $+2.9$. The earlier observation that hyperbolic has \emph{smaller} transfer gap applied to the synthetic-only regime only; once fine-tuned, hyperbolic actually has the larger gap.
\end{itemize}

\paragraph{Why the dim=32 advantage did not compound with fine-tuning.} The synthetic-only result (hyperbolic dim=32 beats Euclidean dim=128 on synthetic) was at least partly an optimisation effect: at dim=128 Euclidean, the model has enough capacity that the optimiser finds a reasonable but not-optimal basin; at dim=32 hyperbolic, the optimiser has less room to wander and lands closer to the geometry's preferred solution. When you introduce fine-tuning, the ``has enough capacity to adapt'' benefit of the larger Euclidean model dominates, and the geometry advantage of hyperbolic becomes secondary.

\paragraph{The defensible claim.} Not ``hyperbolic is better.'' Rather: \textbf{SA-GAT Hyperbolic at output\_dim=32 achieves 0.961 THA on COCO (GT keypoints, after fine-tuning), at $4\times$ smaller embedding than Euclidean dim=128's 0.969}. A $0.8$-point quality cost for a $4\times$ memory/bandwidth saving on the embedding storage. Useful in deployment-constrained settings but not a quality improvement.

\paragraph{Implications for the MEng paper.} The hyperbolic experiments become two ablation rows (synth-only and COCO-ft), framed as an embedding-dimension efficiency study. Not a headline contribution, but a genuine efficiency tradeoff worth documenting. Specifically:
\begin{itemize}
    \item Synthetic-only: hyperbolic dim=32 beats Euclidean dim=128 on every metric. Clean positive result.
    \item COCO-fine-tuned: hyperbolic dim=32 trails Euclidean dim=128 by $\sim 1$ point. Efficiency tradeoff rather than pure improvement.
\end{itemize}

\paragraph{What to do next.} Option A (stop) is now the right call. The hyperbolic experiment is documented, the answer is clear, and pursuing Option C (full manifold architecture) is not justified by these numbers --- if dim=32 hyperbolic with exp-map-only output fine-tunes to 0.961, a full M\"obius-operations architecture would plausibly reach around the same ceiling, at much higher engineering cost. The marginal value is no longer there.

End-to-end with HigherHRNet detection would close the picture ($\sim 0.948$ expected based on the Euclidean $-0.014$ end-to-end drop), but it is not strictly necessary for the ablation-row story.


\subsection{TriGAT --- Triplet-Based Primitive}

\subsubsection{Motivation}

A final architectural experiment from a conversation with a collaborator: instead of embedding joints, embed \emph{triplets} (chains of three skeleton-connected joints A-B-C where A-B and B-C are bones). The angle at the pivot B is rotation-invariant articulation information, which joint embeddings do not directly capture. Two people in a scene with joints in similar absolute positions can still be distinguished by their articulation patterns: different shoulder-elbow-wrist angles, different hip-knee-ankle stances, different head-torso orientations.

The COCO skeleton enumerates exactly 19 canonical triplets (4 head configurations, 4 arm/leg chains, torso configurations around the shoulders and hips). Each person in a scene contributes up to 19 triplet instances (fewer if some joints are invisible). This is a richer primitive than a joint: it encodes geometry and articulation simultaneously.

The pitch was ``we found standard joint embeddings do not capture articulation structure, so we switched to triplet embeddings which do.''

\subsubsection{Implementation}

Created \texttt{trigat.py} with three components:
\begin{itemize}
    \item \texttt{build\_triplet\_graph}: converts a joint PyG graph to a triplet PyG graph. For each person in the scene, enumerates the 19 COCO triplets; skips any with missing joints. Each triplet node has features $[\text{pivot\_x}, \text{pivot\_y}, \Delta_{AB}, \Delta_{CB}, \cos\theta, \sin\theta, |BA|, |BC|]$ = 10 continuous values, plus a learnable 16-dim triplet-type embedding (19 types).
    \item \texttt{TriGATEmbedding}: GATv2-based attention network over the triplet graph. kNN edges over pivot positions, same SA-GAT-like layer stack, outputs per-triplet embeddings.
    \item \texttt{vote\_joint\_labels}: post-clustering step that votes each joint's person assignment from the cluster labels of triplets containing it. Pivot vote weighted $2\times$ over wing votes (pivot's articulation is more distinguishing).
\end{itemize}

Person label for a triplet = pivot joint's person ID. Contrastive loss is the same \texttt{GATOnlyLoss} as for joint embeddings.

Also wrote \texttt{eval\_trigat.py} that runs three grouping methods (kNN, COP-Kmeans, THA) on the triplet embeddings, then votes joint-level labels and computes PGA against the GT joint labels.

\subsubsection{Results}

Trained for 50 epochs on synthetic-only data, same training budget as SA-GAT.

\begin{tabular}{lcc}
\hline
\textbf{Method} & \textbf{Synth PGA} & \textbf{COCO PGA} \\
\hline
\multicolumn{3}{l}{\emph{SA-GAT baselines (synthetic-only, same regime):}} \\
SA-GAT + kNN & 0.910 & 0.882 \\
SA-GAT + COP-Kmeans & 0.925 & 0.915 \\
SA-GAT + THA & 0.938 & — \\
\hline
\multicolumn{3}{l}{\emph{TriGAT (triplet graph + voting):}} \\
Triplets + kNN & 0.898 & 0.855 \\
Triplets + COP-Kmeans & 0.917 & 0.884 \\
Triplets + THA & 0.804 & 0.849 \\
\hline
\end{tabular}

\vspace{0.5em}

\textbf{TriGAT underperforms SA-GAT across every comparison.} The best TriGAT configuration (triplets + COP-Kmeans) scores 0.917 synthetic / 0.884 COCO, compared to SA-GAT's 0.925 / 0.915. The gap widens on COCO, suggesting the triplet representation is losing more information during the sim-to-real transfer than it gains through articulation encoding.

\subsubsection{Why it did not work}

\textbf{1. The voting step is lossy.} Each joint appears in multiple triplets (up to $\sim 3$ triplets per joint). If the triplets for a joint are split across cluster assignments, voting muddies the signal. For a fully-visible person the 19 triplets cover 17 joints unevenly --- the wrists/ankles each appear in 1 triplet, while the shoulders/hips each appear in several. Low-coverage joints depend heavily on a single triplet's cluster assignment.

\textbf{2. THA with voting is catastrophic on synthetic (0.804).} THA's per-type Hungarian matching is decisive: each triplet of type $t$ gets assigned to exactly one cluster. If one triplet in a person is mis-assigned, voting for the joints in that triplet is pulled toward the wrong cluster. On synthetic data with tight, well-separated clusters, the voting amplifies small mistakes rather than absorbing them. COP-Kmeans is more forgiving because its greedy iterative nature produces some softer clusters.

\textbf{3. The angle signal does not transfer well.} The cosine/sine features at the pivot encode local articulation. In synthetic data all people have the same rig and similar pose distributions, so articulation signal is usable. On COCO, real poses vary enormously (sitting, crouching, walking, occluded), and the angle features alone are not a reliable identity marker. The raw position features in SA-GAT (augmented by the attention layers attending over skeletal neighbours) capture absolute spatial structure that transfers more cleanly.

\textbf{4. The triplet type embedding may not be learning what we hoped.} With only 400 synthetic training scenes, the 19-dim triplet-type embedding has ~80 params (19 types $\times$ 16 dims). Each type appears $K$ times per scene, so across 400 scenes with avg 3 people = 1200 examples per type. That should be enough, but the embedding is competing with the raw positional features for gradient signal, and the type embeddings might just be specialising to synthetic-specific angle distributions.

\textbf{5. The graph structure may be fighting itself.} The triplet graph has roughly $19K$ nodes for a $K$-person scene, vs the joint graph's $17K$ nodes. Scale is similar. But the edges are kNN over pivot positions, which means triplets whose pivots are close get attention. In a scene where two people overlap spatially, their triplets' pivots are close but their articulation is different --- exactly the case where triplet embeddings should help. The results suggest that in these overlapping cases, the attention is still dominated by spatial proximity rather than articulation similarity, because the input features are still dominated by positions.

\subsubsection{Analysis}

\textbf{1. This is a clean negative result.} Every TriGAT configuration loses to SA-GAT under identical training conditions. Not a ``close'' experiment --- the gap is consistent and meaningful (1--3 points on strong grouping methods, $>10$ points on THA).

\textbf{2. The architectural idea is not wrong in principle --- the execution has issues.} Angle-based features \emph{should} help for articulated pose, and embedding triplets rather than joints \emph{should} be a richer primitive. But the voting step, the dominance of position features at the input, and the fixed 19-triplet topology (doesn't scale to occluded scenes well) together give up more than the articulation features contribute.

\textbf{3. The single largest problem is almost certainly the voting step.} Triplet-level clustering is not what we want --- we want joint-level assignment. A version that learned the joint assignment directly (e.g., triplet embeddings $\to$ attention $\to$ joint logits) would avoid the vote. But that is essentially reinventing SA-GAT with a different parametrisation, and the more parsimonious conclusion is that joint embeddings are the right primitive.

\textbf{4. The experiment also informs the paper-writing approach.} We have now tested several alternative embedding primitives:
\begin{itemize}
    \item Joint position + type (SA-GAT) --- baseline, works well
    \item Joint + visual features (SA-GAT v2) --- marginal gain, does not beat the simple version
    \item Lorentz-hyperbolic joint (SA-GAT hyperbolic) --- efficiency trade, not a quality win
    \item Triplet (TriGAT) --- clean underperformance
\end{itemize}
This is useful ablation evidence in the MEng paper: ``we investigated four alternative embedding formulations and the simplest one (positional joint embeddings) remained strongest.'' That is a valid empirical claim.

\textbf{5. Recommendation: close this direction.} Documented as a negative result in the dissertation. The MEng paper can mention it in one sentence during the ablation section (``triplet-based embeddings with majority-vote grouping underperform joint embeddings by $\sim 1$--$3$ points across all methods tested''). No follow-up experiments justified.


\subsection{Publication Plan}

The publication plan mirrors the MEng/PEng conceptual split. The MEng paper presents pose grouping as an embedding problem and shows SA-GAT is the right embedding network for it. The PEng paper reframes the problem as structured assignment and presents THA as the right algorithm for that framing, with SCOT and the investigation behind it as the experimental evidence that motivates the reframing.

\subsubsection{MEng paper: SA-GAT (target: submitted)}

\textbf{Working title:} ``SA-GAT: Skeleton-Aware Graph Attention Networks for Multi-Person Pose Grouping''

\textbf{Core argument:} Multi-person pose grouping is structurally an embedding problem. If the embeddings are good enough, simple clustering outperforms every learned grouping head. SA-GAT is a GATv2-based embedding network with three lightweight skeleton-aware modifications that produce embeddings of this quality, and the resulting SA-GAT + constrained clustering pipeline matches state-of-the-art end-to-end methods while remaining fully decoupled from the keypoint detector.

\textbf{Narrative arc:}
\begin{enumerate}
    \item \textbf{Motivation.} Bottom-up pose grouping has traditionally been studied by designing complex grouping heads (slot attention, associative embedding, modularity clustering, edge classification). The paper argues that the embedding space is the real lever.
    \item \textbf{Baseline and failed alternatives.} Standard GAT + contrastive loss + kNN clustering gives a strong baseline (0.879 COCO). Learned grouping heads on top (slot attention, graph partitioning, DEC) all fail to improve on kNN. SA-DMoN with modularity-based clustering fails across 10 experimental variants because modularity is misaligned with identity-based grouping.
    \item \textbf{The insight.} When the embeddings are well-structured by contrastive learning, simple constrained clustering (COP-Kmeans with must-not-link constraints between same-type keypoints) outperforms every learned head tested.
    \item \textbf{SA-GAT.} Three skeleton-aware modifications (type-pair attention bias, skeleton-relative position encoding, same-type repulsion heads) each contribute independently to embedding quality. Total cost: $\sim$10K extra parameters over standard GAT.
    \item \textbf{Synthetic-to-COCO transfer.} Depth removal fixed the MiDaS domain gap (+4.1 points). COCO fine-tuning closed the transfer gap almost entirely, reaching 0.968 COCO PGA with COP-Kmeans.
    \item \textbf{End-to-end evaluation.} Paired with HigherHRNet-detected keypoints, SA-GAT + COP-Kmeans achieves 0.955 end-to-end PGA, within 3 points of HigherHRNet's own associative embedding grouping (0.985).
    \item \textbf{Modularity claim.} SA-GAT is a detector-agnostic grouping module, in contrast to jointly-trained AE heads that are locked to their specific backbone.
\end{enumerate}

\textbf{Key results:}
\begin{itemize}
    \item 0.968 COCO PGA with GT keypoints (COP-Kmeans on COCO-fine-tuned SA-GAT)
    \item 0.955 COCO PGA end-to-end with HigherHRNet detection
    \item SA-GAT ablation: each of the three modifications contributes
    \item Every learned grouping head tested fails to beat simple clustering on the same embeddings
    \item Detector-agnostic: trained independently of the keypoint detector, works as a drop-in module
\end{itemize}

\textbf{Target venue:} \textbf{Image and Vision Computing} (Elsevier). Q1 in Computer Vision and Pattern Recognition (SCImago 2024), IF 5.34, SJR 0.791, h-index 150. Scopus-indexed, DHET-accredited. Publishes pose estimation work regularly and fits the architectural-contribution-plus-empirical-study format. Pattern Recognition (Q1, IF 9.84) is the stretch target.

\textbf{Key strengths for review:}
\begin{itemize}
    \item A coherent argument supported by negative results (failed learned heads, failed modularity)
    \item Novel architecture with clean ablation
    \item Strong empirical results on COCO (both GT keypoints and end-to-end with an external detector)
    \item Lightweight and modular, unlike jointly-trained end-to-end methods
\end{itemize}


\subsubsection{PEng paper: SCOT \& THA (target: published)}

\textbf{Working title:} ``From Clustering to Assignment: Skeleton-Constrained Optimal Transport and Type-wise Hungarian Matching for Multi-Person Pose Grouping''

\textbf{Core argument:} Multi-person pose grouping is not a clustering problem --- it is a structured assignment problem where each person has exactly 17 typed slots and the task is to assign detected keypoints to slots under hard type constraints. This reframing motivates a family of assignment-based algorithms. SCOT formalises the framing as skeleton-constrained optimal transport; it is novel but underperforms due to specific architectural limitations. Diagnosing those limitations through systematic ablation leads to THA, a type-decomposed Hungarian matching algorithm that realises the assignment-based view of the problem and outperforms clustering-based methods across every condition tested.

\textbf{Narrative arc:}
\begin{enumerate}
    \item \textbf{The reframing.} Hard type constraints (no two left elbows in the same person) are not a clustering property, they are an assignment property. Pose grouping is structurally closer to linear assignment than to k-means. This reframing is what motivates the whole paper.
    \item \textbf{SCOT: formalising the assignment view.} A differentiable optimal transport formulation with typed skeleton slots ($K \times 17$), infinite cost masking for type violations, and Sinkhorn iterations. Novel: no prior work applies structured OT with typed slots to keypoint grouping. Good synthetic performance (0.923 PGA).
    \item \textbf{SCOT fails to transfer.} On COCO, SCOT drops to 0.862 --- far below the clustering-based COP-Kmeans baseline (0.910). This is the puzzle the paper investigates.
    \item \textbf{Systematic diagnosis.} Seven SCOT variants (residual, adaptive, unbalanced, dustbin, k-means initialised), three training strategies (joint synthetic, joint COCO fine-tune, frozen backbone), two backbones (SA-GAT, standard GAT), two embedding sources (our contrastive embeddings, HigherHRNet's own backbone features). All land in the same 0.84--0.88 COCO range. The bottleneck is architectural, not training or embedding.
    \item \textbf{Identified failure modes.} (a) Fixed learned prototypes cannot represent the diversity of real scenes. (b) Balanced Sinkhorn transport assumes uniform joint visibility, which is violated under occlusion. (c) Soft-to-hard argmax conversion loses information near decision boundaries.
    \item \textbf{SCOT-KI as diagnostic step.} Replacing fixed prototypes with per-scene k-means centroids recovers 6.5 points (0.861 $\to$ 0.926). Confirms prototype rigidity was the first problem; remaining gap to COP-Kmeans is in the Sinkhorn mechanism itself.
    \item \textbf{THA: the assignment-based algorithm that works.} Per-joint-type bipartite matching via Hungarian, with shared scene-adaptive centroids and iterative refinement. Scene-adaptive (inheriting from COP-Kmeans), globally optimal per-type (inheriting the assignment framing from SCOT), and structurally constrained (each type matched independently). Zero learned parameters.
    \item \textbf{Novelty verification.} THA's combination of per-type decomposition, Hungarian optimal assignment, and shared iterative centroids has not been published. Related methods (OpenPose pairwise matching, COP-Kmeans greedy constraint resolution, DETR-style Hungarian for set prediction, DC-GNN, HGG) are structurally different.
    \item \textbf{Results.} THA beats COP-Kmeans on synthetic (+2.7), matches on COCO with GT keypoints (+0.1), matches end-to-end ($\approx$0.955).
\end{enumerate}

\textbf{Key results:}
\begin{itemize}
    \item SCOT: novel OT formulation for pose grouping, 0.923 synthetic / 0.862 COCO (GT keypoints)
    \item SCOT-KI: 0.926 COCO (scene-adaptive prototypes recover 6.5 points)
    \item THA: 0.965 synthetic / 0.969 COCO / 0.955 end-to-end
    \item Bulletproof SCOT diagnosis: same ranking (SCOT-KI $<$ THA $\approx$ COP-Kmeans) across 7 variants, 3 training strategies, 2 backbones, 2 embedding sources
    \item SA-DMoN investigation (10 experiments) as additional evidence that clustering-based approaches hit a ceiling --- included as supporting material
\end{itemize}

\textbf{Target venue:} \textbf{Pattern Recognition Letters} (Elsevier). Q1 in Computer Vision and Pattern Recognition (SCImago 2024), IF 4.54, SJR 1.005, h-index 188. Scopus-indexed, DHET-accredited. The shorter format suits a paper built around a conceptual reframing with experimental evidence rather than a large system. Neurocomputing (Q1, IF 8.18) is an alternative if a longer format is preferred.

\textbf{Key strengths for review:}
\begin{itemize}
    \item A clear conceptual contribution (pose grouping as structured assignment, not clustering) supported by both theory and experiment
    \item Two novel algorithms (SCOT, THA) with different roles: SCOT establishes the framing, THA realises it in a form that actually works
    \item Exceptionally thorough experimental methodology --- 7 SCOT variants, 3 training strategies, 2 backbones, 2 embedding sources
    \item Honest reporting: the paper does not hide SCOT's underperformance, it uses it as diagnostic evidence for the architectural insights that produce THA
    \item THA beats the best clustering baseline (COP-Kmeans from the MEng paper), so the conceptual reframing has measurable payoff
\end{itemize}


\subsubsection{Relationship between the papers}

The two papers mirror the MEng/PEng conceptual split.

The MEng paper (SA-GAT) adopts the standard clustering framing of pose grouping and contributes a better embedding network. It is self-contained: SA-GAT + COP-Kmeans is a complete, empirically strong system that does not depend on anything from the PEng paper.

The PEng paper (SCOT/THA) takes SA-GAT as a given (reference is made to the MEng paper) and argues that the clustering framing misses what the problem actually is. The paper reframes pose grouping as structured assignment, presents SCOT as the formal OT realisation of that framing, diagnoses why SCOT does not work as-is, and presents THA as the algorithm that actually realises the assignment-based framing in a form that outperforms clustering baselines.

The papers share infrastructure (synthetic data generator, COCO adapter, evaluation protocol, end-to-end pipeline) but have independent contributions and can each be read standalone. Read together, they tell the full story: \emph{produce good embeddings (MEng), then reason about how to assign them to people (PEng).}

For the PEng degree specifically, the paper-1 / paper-2 structure satisfies the ``one published, one submitted'' requirement while giving each paper a clear and different contribution --- the MEng paper is system-focused (architecture + empirical results), the PEng paper is conceptually focused (reframing + algorithm + diagnostic ablation).

\subsubsection{Venue summary}

All target journals are Elsevier, Scopus-indexed, and on the DHET accredited list (verified via the ISI/Scopus sublists that DHET uses for subsidy eligibility). Quartile rankings are from SCImago 2024.

\begin{tabular}{lllll}
\hline
\textbf{Journal} & \textbf{Quartile} & \textbf{IF} & \textbf{SJR} & \textbf{Target for} \\
\hline
Image and Vision Computing & Q1 & 5.34 & 0.791 & Paper 1 (primary) \\
Pattern Recognition & Q1 & 9.84 & 2.058 & Paper 1 (stretch) \\
Pattern Recognition Letters & Q1 & 4.54 & 1.005 & Paper 2 (primary) \\
Neurocomputing & Q1 & 8.18 & 1.471 & Paper 2 (alternative) \\
\hline
\end{tabular}

\vspace{0.5em}

All four are Q1. The primary targets (Image and Vision Computing, Pattern Recognition Letters) are realistic given the results and novelty. The stretch/alternative targets (Pattern Recognition, Neurocomputing) have higher impact factors but are more competitive.


\subsubsection{Naming concern --- is ``SA-GAT'' too generic?}

Before submitting the MEng paper I need to consider whether the name ``SA-GAT'' is distinctive enough. ``Skeleton-aware graph attention'' is a general concept that several existing works have already explored in adjacent settings.

\paragraph{Existing adjacent work.} Notable collisions in the literature:
\begin{itemize}
    \item \textbf{STGAT / ST-GAT} (Hu et al., \emph{Expert Systems with Applications} 2023)~\cite{hu2023stgat} --- spatial-temporal graph attention for skeleton-based action recognition. Attends over the same COCO-style skeleton graph as SA-GAT but for a completely different task (action classification from joint trajectories).
    \item \textbf{GAST-Net} (Liu et al., arXiv 2020)~\cite{liu2020gastnet} --- graph attention spatio-temporal convolutional network for single-person 3D pose estimation from video. Uses graph attention to model kinematic constraints within a single person's skeleton.
    \item \textbf{DGAN, MGSAN, and similar} --- various graph attention variants for skeleton-based action recognition with names that partially overlap.
\end{itemize}

The broader concept of attending over skeleton joints via graph attention is firmly established prior art. A casual reader searching for ``skeleton-aware graph attention'' would find these papers first.

\paragraph{What is genuinely novel in our SA-GAT.}
\begin{itemize}
    \item The three specific architectural modifications --- type-pair attention bias, skeleton-relative position encoding, same-type repulsion heads --- as a package.
    \item The \emph{repulsion heads} idea specifically (pushing same-type keypoints apart for identity separation) is different from existing methods that attend along skeletal connections.
    \item The \emph{application} to bottom-up multi-person person-grouping. None of the cited works address grouping; they are all single-person or action recognition.
\end{itemize}

\paragraph{What is not novel.}
\begin{itemize}
    \item Graph attention over skeleton joints (many precedents).
    \item The phrase ``skeleton-aware'' as a descriptor (widely used).
\end{itemize}

\paragraph{Options to address this.}

\emph{Option A: Rename to something task-specific.} Candidate names that avoid existing collisions and signal the application clearly:
\begin{itemize}
    \item \textbf{PG-GAT} --- Pose Grouping GAT. Most literal, short enough for tables.
    \item \textbf{GroupGAT} --- emphasises the grouping task.
    \item \textbf{ID-GAT} --- Identity GAT, signalling the identity-separation objective.
    \item \textbf{AssignGAT} --- would pair well with the PEng paper's assignment framing.
\end{itemize}

\emph{Option B: Keep SA-GAT but be explicit in the paper.} Retain the name, but add a paragraph early in the paper that acknowledges the prior art and specifies our contribution as the application and the three specific modifications rather than the general mechanism. Example phrasing: ``We retain `skeleton-aware graph attention' as a descriptor of the mechanism, but our contribution is not the attention mechanism itself (which has antecedents in STGAT~\cite{hu2023stgat} for action recognition and GAST-Net~\cite{liu2020gastnet} for single-person pose) but rather its adaptation to bottom-up multi-person keypoint grouping through three specific modifications.''

\paragraph{Recommendation.} Lean toward Option A, specifically \textbf{PG-GAT} or \textbf{GroupGAT}. Naming collisions hurt discoverability --- if the paper gets indexed alongside action-recognition STGAT variants it will be harder to find via keyword search. A task-specific rename is a cheap change (a few hours of codebase rename + consistent terminology in the paper) that could meaningfully help the paper stand out.

If we do rename, the search-and-replace is straightforward because the module is well-isolated: \texttt{sa\_gat.py}, \texttt{sa\_gat\_v2.py}, the config classes \texttt{SAGATConfig} and \texttt{SAGATV2Config}, and the class \texttt{SAGATEmbedding}. The output directory structure stays the same. No experimental results need to be re-run.

The final decision can wait until we start writing the MEng paper, but the question is flagged here so it is not forgotten.


\subsection{COCO Evaluation Composition --- Single-Person Images}

\subsubsection{Motivation}

A code review of the evaluation pipeline identified an important detail in the COCO protocol.
The COCO validation evaluation currently includes all valid person-keypoint images, including images with only one annotated person.
This is because \texttt{CocoAdapter} defaults to \texttt{min\_people=1}, and the main evaluation scripts instantiate it without overriding that value.

This is not a bug in the sense of an incorrect metric implementation.
All methods reported so far were evaluated under the same protocol, so the comparisons between methods remain internally fair.
However, the detail matters for interpretation because the dissertation is about \emph{multi-person pose grouping}.
When an image contains only one person, the grouping problem is almost trivial: every keypoint belongs to the same person.
There is no same-type conflict, no overlapping-person ambiguity, and no meaningful assignment decision between people.

The consequence is that COCO-all PGA is an easier metric than a strict multi-person grouping metric.
It measures performance over the natural filtered COCO validation distribution used by the project, but it does not isolate the hard grouping cases.

\subsubsection{Dataset composition}

The COCO val2017 person-keypoint annotations contain a large fraction of single-person images under the current filtering.
Counting valid person-keypoint annotations gives:

\begin{tabular}{lcccc}
\hline
\textbf{Filter} & \textbf{Images with $\geq 1$ person} & \textbf{Single-person} & \textbf{Multi-person} & \textbf{Single-person fraction} \\
\hline
Any keypoint annotation & 2346 & 1098 & 1248 & 46.8\% \\
At least 3 labelled keypoints & 2258 & 1063 & 1195 & 47.1\% \\
\hline
\end{tabular}

\vspace{0.5em}

Thus, almost half of the current COCO-all evaluation set consists of images that do not stress the grouping component.
This explains why COCO-all scores are noticeably higher than scores on the multi-person-only subset.

\subsubsection{Diagnostic results}

The existing COCO-all headline was compared with a stricter \texttt{min\_people=2} evaluation on COCO val2017.
The \texttt{min\_people=2} setting retains 1248 images and removes the single-person cases.
The hyperbolic rows use the raw tangent-space diagnostic from the code review, which preserves the radial information discarded by the original L2-normalised tangent evaluation.

\begin{tabular}{lcccc}
\hline
\textbf{Model / grouping method} & \textbf{kNN} & \textbf{COP-KM} & \textbf{THA} & \textbf{Subset} \\
\hline
Euclidean PG-GAT, COCO-ft & 0.957 & 0.968 & 0.969 & $\geq 1$ person \\
Euclidean PG-GAT, COCO-ft & 0.921 & 0.941 & 0.943 & $\geq 2$ people \\
\hline
Hyperbolic PG-GAT, COCO-ft (raw tangent) & 0.947 & 0.962 & 0.965 & $\geq 1$ person \\
Hyperbolic PG-GAT, COCO-ft (raw tangent) & 0.902 & 0.929 & 0.935 & $\geq 2$ people \\
\hline
\end{tabular}

\vspace{0.5em}

The absolute numbers drop on the multi-person-only subset, as expected.
For the Euclidean COCO-fine-tuned PG-GAT model, THA drops from 0.969 to 0.943.
For the hyperbolic COCO-fine-tuned PG-GAT model, THA drops from 0.965 to 0.935.
The ranking remains stable: Euclidean PG-GAT remains stronger than the hyperbolic variant after COCO fine-tuning, and THA remains stronger than kNN on the stricter subset.

\subsubsection{Interpretation}

\textbf{1. COCO-all is still a reasonable benchmark.} COCO val2017 is a standard evaluation source, and retaining all valid person-keypoint images gives a reproducible distribution. The existing headline numbers are therefore not invalid. They should simply be described as COCO-all results rather than as a pure multi-person grouping stress test.

\textbf{2. Single-person images inflate grouping PGA.} The inflation is structural, not accidental. A single-person scene has $K=1$, so grouping reduces to assigning all nodes to one cluster. This makes the metric easier because the hard part of bottom-up grouping --- separating keypoints from different people --- is absent.

\textbf{3. The stricter subset is the better grouping diagnostic.} The \texttt{min\_people=2} subset directly tests the multi-person grouping problem. It should be reported as a secondary table in the dissertation and paper, especially in the results or discussion section, to show how the method behaves when the grouping ambiguity is actually present.

\textbf{4. The main conclusions survive.} The stricter subset lowers the absolute scores but does not reverse the important ordering. PG-GAT with THA remains better than PG-GAT with kNN, and the hyperbolic variant remains an efficiency trade rather than a quality win. The conclusion becomes more conservative, not weaker.

\textbf{5. CrowdPose is not necessary for the MEng dissertation.} CrowdPose would be a harder crowd-focused benchmark, but adding it would expand the dataset pipeline and evaluation scope late in the project. The cleaner MEng solution is to keep COCO-all for comparability and add COCO multi-person-only as a diagnostic. CrowdPose can be listed as future work or reserved for a later PEng/PhD extension.

\subsubsection{Recommendation for dissertation reporting}

The dissertation should report the COCO results in two layers:

\begin{itemize}
    \item \textbf{COCO-all:} the existing headline protocol, retained for continuity with the experimental record and comparison with HigherHRNet AE on the same distribution.
    \item \textbf{COCO multi-person-only:} a stricter diagnostic table using \texttt{min\_people=2}, used to isolate the true multi-person grouping difficulty.
\end{itemize}

Suggested wording for the dissertation:

\begin{quote}
COCO val2017 contains both single-person and multi-person images.
The main COCO results retain all valid person-keypoint images for comparability with the detector-based evaluation.
Because single-person images do not require a non-trivial grouping decision, a second evaluation is reported on the subset with at least two annotated people.
This subset provides a stricter diagnostic of the multi-person grouping problem.
\end{quote}

This avoids overstating the COCO-all numbers while preserving the existing experimental record.

\subsection{Weights \& Biases Integration and Hyperparameter Sweeps --- Towards Defensible Configuration Choices}

\subsubsection{Motivation}

Up to this point I have run every experiment by hand: write a YAML config, run \texttt{trainer.py}, look at the JSON output, eyeball whether it beat the previous best, and update the lab book. Two problems are now blocking dissertation writeup:

\textbf{1. ``Best'' and ``latest'' are not the same thing.} The current output structure uses a \texttt{latest} symlink that always points to the most recent run of a given config. If I retrain after tweaking something and the run is worse, \texttt{latest} silently demotes the better checkpoint. I have already lost results this way once (the original SCOT 0.923 was overwritten before I introduced the GUID-based output system). I need a separation of ``most recent run of this config'' from ``best run for this config'' that does not rely on me remembering to copy files into \texttt{frozen\_checkpoints/}.

\textbf{2. I cannot defend the hyperparameter choices.} For every architectural decision in PG-GAT (depth, heads, hidden dim, output dim, contrastive temperature, type-pair coefficient, repulsion weight, fine-tuning learning rate) the honest answer to ``why did you pick that value'' is currently ``it worked.'' That is not a defensible answer at the dissertation oral and it is not a publishable one. A Bayesian sweep over each axis, logged centrally, gives me a parallel-coordinate plot, a parameter-importance ranking, and the headline ``best run achieved $X$ at config $Y$.''

\textbf{3. Cross-experiment comparison is hand-managed.} Comparing PG-GAT + \{kNN, COP-Kmeans, THA\} across the same checkpoint currently means opening three JSON files. With many checkpoints this is error-prone.

Weights \& Biases (\texttt{wandb}) addresses all three: artifact aliases (\texttt{best}, \texttt{latest}) replace the symlink convention; sweep configs replace manual iteration; the dashboard replaces JSON file diving. Free academic tier covers everything I need.

\subsubsection{Integration plan}

Minimum-viable diff to the existing codebase, preserving backwards compatibility with all 42 existing YAML configs:

\begin{itemize}
    \item Add \texttt{wandb: Optional[WandbConfig] = None} to \texttt{ExperimentConfig} in \texttt{config.py}.
    \item Guard all wandb calls in \texttt{trainer.py} with \texttt{if cfg.wandb is not None}, so existing configs still run unchanged.
    \item Initialise with \texttt{wandb.init(project="...", name=guid, config=cfg.model\_dump(), save\_code=True)}; the GUID stays the run identity.
    \item Log per-epoch metrics: \texttt{train/loss}, \texttt{val/PGA\_synth}, \texttt{val/PGA\_coco}, and the K-stratified \texttt{val/PGA\_K=1}, \texttt{val/PGA\_K=2}, \texttt{val/PGA\_K=3}, \texttt{val/PGA\_K\_ge\_4}.
    \item Log checkpoints as artifacts on validation improvement, with alias \texttt{latest} auto-rotating; \texttt{best} alias set manually after a sweep finishes.
    \item GUID directory structure on disk is unchanged; W\&B is a remote indexed layer over local files. If wandb sync fails, training continues.
\end{itemize}

The brief from a separate planning agent proposed a more aggressive integration including dataset versioning and a full evaluation refactor. I am skipping the dataset versioning (synthetic generator is in git, COCO is canonical) and keeping the existing 10 \texttt{eval\_*.py} scripts as-is, with a thin \texttt{eval\_by\_run.py} wrapper that resolves a run reference (alias, run ID, or GUID) to a checkpoint path and dispatches.

\subsubsection{Project, group, and tag conventions}

\begin{tabular}{ll}
\hline
\textbf{Field} & \textbf{Value} \\
\hline
Project & \texttt{multiperson-pose-grouping} \\
Groups & \texttt{pg-gat-arch-sweep}, \texttt{pg-gat-loss-sweep}, \texttt{pg-gat-finetune-sweep},\\
& \texttt{scot-diagnosis}, \texttt{tha-eval}, \texttt{hyperbolic-ablation}, \texttt{trigat-ablation} \\
Tags & \texttt{meng-paper}, \texttt{peng-paper}, \texttt{dissertation-final}, \\
& \texttt{paper-final-meng}, \texttt{paper-final-peng}, \texttt{ablation}, \texttt{negative-result} \\
Aliases & \texttt{latest} (auto), \texttt{best} (after sweep), \texttt{champion}, \\
& \texttt{meng-headline}, \texttt{peng-headline} \\
Sweep metric & \texttt{val/PGA\_coco} (primary); K-stratified synth as secondary \\
\hline
\end{tabular}

\vspace{0.5em}

I picked \texttt{val/PGA\_coco} as the sweep selection metric over the brief's proposed \texttt{val/PGA\_K\_ge\_2} (synth) because the dissertation headline is the COCO number. Selecting on synth and shipping to COCO is what produced the depth-channel surprise (synth said yes, COCO said no, $-4.1$ points). Selecting directly on COCO is honest and matches the dissertation's primary metric. K-stratified synth PGA is still logged so I can show robustness across difficulty.

I picked one project rather than splitting MEng and PEng because the work is one body of research. The PEng pivot (COP-Kmeans, SCOT, THA) explicitly builds on MEng PG-GAT embeddings, and the cleanest cross-cuts (THA on hyperbolic embeddings, SCOT on HRNet features) straddle both papers. Tag separation handles paper boundaries; project separation would force duplicate logging.

\subsubsection{Planned sweeps}

\begin{tabular}{lll}
\hline
\textbf{Sweep} & \textbf{Axes} & \textbf{Defends} \\
\hline
PG-GAT architecture & depth $\in \{2,3,4,5\}$, heads $\in \{2,4,8\}$, & RQ2 (which mods help, by how much), \\
                    & hidden dim $\in \{64,128,256\}$, output dim $\in \{32,64,128\}$ & basic architecture defence \\
PG-GAT loss        & contrastive temp $\in \{0.1, 0.2, 0.5\}$, & loss-coefficient choices in \\
                    & repulsion weight $\in \{0.1, 0.5, 1.0, 2.0\}$, & Ch3 PG-GAT writeup \\
                    & type-pair coefficient $\in \{0.1, 0.5, 1.0, 2.0\}$ & \\
COCO fine-tuning   & epochs $\in \{10,20,30,50\}$, & RQ6 (how much does fine-tune lift, \\
                    & lr $\in [10^{-5}, 10^{-3}]$ log-uniform, & does it damage synth) \\
                    & freeze-strategy $\in \{$none, backbone, encoder$\}$ & \\
Hyperbolic dim     & dim $\in \{16, 32, 64, 128\}$, both pre-ft and post-ft & RQ4 (hyperbolic vs Euclidean,  \\
                    & & efficiency vs quality)\\
\hline
\end{tabular}

\vspace{0.5em}

All sweeps use Bayesian search with hyperband early-termination. Each sweep budget capped at 32 runs (architecture, loss) or 16 runs (fine-tune, hyperbolic) to keep total compute bounded. Estimate: $\approx$~32 + 32 + 16 + 16 = 96 GPU-hours across all four sweeps at $\sim$30 minutes per run on the 5070 Ti.

\subsubsection{Defensibility argument}

\textbf{1. Each sweep has a research question it defends.} Every choice the dissertation makes about PG-GAT can be traced to one of the four sweeps above. ``Why depth=4'' becomes ``Bayesian sweep over $\{2,3,4,5\}$, winning depth was 4, parallel-coordinate plot in Appendix B.'' The four sweeps cover RQ2, RQ4, and RQ6 directly; RQ1 (embedding vs head) is already settled by the failed-heads chapter; RQ3 (detector independence) and RQ5 (graph primitive) are settled by the E2E and TriGAT experiments respectively.

\textbf{2. The selection metric matches the headline metric.} Sweeping on \texttt{val/PGA\_coco} means the winning configuration is the one that scored best on the same metric the dissertation uses to claim its headline result. This eliminates the synth-versus-COCO mismatch that bit me on the depth channel.

\textbf{3. The sweeps are recoverable post-hoc.} Each sweep run's full config is stored in W\&B. Six months from now, when a reviewer asks ``did you try $X$,'' I can answer either ``yes, run \texttt{abc123}, scored $Y$'' or ``no, but the parameter-importance plot suggests it would not have helped because $Z$.'' Without W\&B that answer is not retrievable.

\textbf{4. The export is appendix-ready.} After all sweeps finish, \texttt{export\_results.py} dumps every \texttt{paper-final}-tagged run to a CSV with full hyperparameters and metrics. That CSV is Appendix B verbatim.

\subsubsection{What to do next}

\begin{itemize}
    \item Create the W\&B account and the \texttt{multiperson-pose-grouping} project.
    \item Write the trainer instrumentation and the three sweep YAMLs (architecture, loss, fine-tune). Hyperbolic sweep is lower priority --- the existing single dim=32 run is already a written ablation.
    \item Run the architecture sweep first --- it directly constrains every later experiment.
    \item Promote the winning configurations as \texttt{champion} and \texttt{meng-headline} aliases; rerun all evaluation scripts (\texttt{evaluator.py}, \texttt{eval\_cop\_kmeans.py}, \texttt{eval\_hungarian\_grouping.py}, \texttt{eval\_end\_to\_end.py}) on the champion checkpoint.
    \item Write Ch3 PG-GAT and Ch4 ablations against the sweep results, citing W\&B run IDs throughout.
\end{itemize}

The existing \texttt{outputs/frozen\_checkpoints/} directory remains the on-disk source of truth; W\&B is the remote, indexed, queryable copy. Belt and braces.

\subsubsection{Implementation update (2026-05-05)}

I wired the W\&B integration end-to-end. The following files were added or modified:

\begin{tabular}{ll}
\hline
\textbf{File} & \textbf{Change} \\
\hline
\texttt{code\_v3/config.py} & Added \texttt{WandbConfig}; \texttt{wandb: Optional[WandbConfig] = None} on \texttt{ExperimentConfig} \\
\texttt{code\_v3/trainer.py} & \texttt{wandb.init} guarded by \texttt{if cfg.wandb is not None}; per-epoch \\
& \texttt{train/loss}, \texttt{train/gat\_loss}, \texttt{train/head\_loss}, \texttt{lr}; \texttt{val/pga} on \\
& val epochs; checkpoint artifact upload with \texttt{latest} alias on each \\
& new best; \texttt{wandb.finish()} with summary fields. New \texttt{--set} flag \\
& applies dotted-path overrides for sweep agents. \\
\texttt{code\_v3/scripts/promote\_best.py} & New. Walks all versions of an artifact (or all model artifacts in a \\
& sweep), picks the highest \texttt{val\_pga} from metadata, applies the \\
& \texttt{best} (or named) alias. \\
\texttt{code\_v3/configs/train\_sa\_gat\_sweep\_base.yaml} & New. Base config that sweeps build on (PG-GAT full + wandb section). \\
\texttt{code\_v3/configs/sweeps/pg\_gat\_arch.yaml} & New. Bayes sweep over depth, heads, hidden, output, joint emb, dropout. \\
\texttt{code\_v3/configs/sweeps/pg\_gat\_loss.yaml} & New. Bayes sweep over contrastive margin, repulsion heads, lr, weight decay. \\
\texttt{code\_v3/configs/sweeps/pg\_gat\_finetune.yaml} & New. Bayes sweep over fine-tune epochs, lr, weight decay; loads \\
& \texttt{frozen\_checkpoints/sa\_gat\_full/best.pt} as pretrained. \\
\texttt{code\_v3/configs/train\_sa\_gat\_smoke\_wandb.yaml} & New. 5-epoch smoke test with wandb enabled. \\
\texttt{.gitignore} & Added \texttt{wandb/} and \texttt{code\_v3/wandb/}. \\
\hline
\end{tabular}

\vspace{0.5em}

\textbf{1. Backwards compatibility verified.} All 42 existing YAML configs still load with \texttt{cfg.wandb = None}, which keeps every wandb branch in \texttt{trainer.py} dormant. The check was a small Python script that loaded each YAML and asserted \texttt{cfg.wandb is None}.

\textbf{2. Smoke test passed.} Ran \texttt{trainer.py --config configs/train\_sa\_gat\_smoke\_wandb.yaml} (5 epochs, val every epoch, PG-GAT full). The W\&B dashboard showed: the run under group \texttt{smoke-test} with the GUID as the run name; per-epoch loss curves; the full Pydantic config in the config tab; three artifact versions (\texttt{v0}, \texttt{v1}, \texttt{v2}) corresponding to the three improvement events during the run; \texttt{latest} alias on \texttt{v2}; clean finish with \texttt{best\_val\_pga} in the run summary.

\textbf{3. Override mechanism works.} The \texttt{--set} flag applies dotted-path overrides to the loaded config dict before Pydantic instantiation. Verified that \texttt{--set sa\_gat.num\_layers=4 --set training.lr=1e-3} correctly modifies the config; the coercion rule converts strings to int, float, bool, or null where appropriate. This is the entry point for wandb sweep agents --- the sweep YAMLs use \texttt{command:} blocks that expand each parameter as a \texttt{--set} argument.

\textbf{4. ``best'' versus ``latest'' clarified during implementation.} Within a single training run \texttt{latest} \emph{is} the best of that run, because \texttt{best.pt} is only overwritten when validation improves and we only upload to W\&B on improvement. The cross-run case is the only place where the two diverge: if a retrain produces a worse checkpoint, \texttt{latest} updates and \texttt{best} stays put. \texttt{promote\_best.py} closes that gap by selecting the version with the highest \texttt{val\_pga} in metadata across all versions of an artifact (or all model artifacts in a sweep).

\textbf{5. Selection metric.} The sweeps select on \texttt{val/pga}, which under the current trainer is the synth val PGA. The COCO and end-to-end numbers are produced by re-running \texttt{eval\_cop\_kmeans.py}, \texttt{eval\_hungarian\_grouping.py}, and \texttt{eval\_end\_to\_end.py} on the winning checkpoint after the sweep finishes. Synth val is fast enough to run every two epochs without dominating sweep time, but it is not the dissertation headline --- the headline is generated by post-hoc evaluation of the winning checkpoint on COCO val2017. This is the same protocol that produced the existing \texttt{train\_sa\_gat\_full} numbers; the sweep does not change which numbers go into the dissertation, only how the winning configuration is selected.

\textbf{6. Known limitation: list-valued overrides.} The current \texttt{\_coerce} helper does not parse list literals such as \texttt{["sweep","arch","meng-paper"]}, so \texttt{--set wandb.tags=...} would fail Pydantic validation. The three sweep YAMLs currently rely on \texttt{wandb.group} for organisation (one group per sweep) and skip per-run tags. This is sufficient because the W\&B UI filters cleanly by group, but it should be fixed before the dissertation-final tagging pass.

\subsubsection{Status snapshot}

\begin{itemize}
    \item Done: account, project, gitignore, trainer instrumentation, sweep base config, three sweep YAMLs, \texttt{promote\_best.py}, smoke test passed.
    \item Pending: launching the architecture sweep (32 runs, $\sim$16 GPU-hours); promoting its winner; updating the loss and fine-tune sweep YAMLs to use the architecture-sweep winning configuration.
    \item Deferred: eval-side wandb hooks (so COCO and end-to-end PGA show up in the same run as synth val); \texttt{eval\_by\_run.py} wrapper (resolve \texttt{--run \{best, id, guid\}} to a checkpoint path); \texttt{scripts/export\_results.py} for the dissertation appendix CSV. These are worth wiring once the first sweep produces a champion to point them at.
\end{itemize}

\subsection{PG-GAT Architecture Sweep --- A Defensible Configuration}

\subsubsection{Motivation}

Every architectural choice in PG-GAT (depth, attention heads, hidden dimension, output dimension, joint embedding dimension, dropout) was originally hand-picked.
The frozen \texttt{train\_sa\_gat\_full} configuration --- 2 layers, 4 heads, hidden 64, output 128, joint embedding 32, dropout 0.1 --- reached 0.910 synth val and 0.882 COCO kNN on the original synthetic-only training run.
That number has anchored every later experiment in this lab book.
But the honest answer to ``why those values'' has always been ``it worked.''
That is not enough for the dissertation defence.
This sweep is the structured search that produces a defendable answer to RQ2 (which architectural choices matter, by how much).

\subsubsection{Setup}

Bayesian sweep with hyperband early termination, 32 runs, selection metric \texttt{val/pga} (synth val, the trainer's primary metric).
Same 50-epoch training schedule as the frozen baseline (\texttt{train\_sa\_gat\_full.yaml}) so the winning configuration is directly comparable to the existing 0.910 reference.
Hyperband \texttt{min\_iter=5} (eligible for termination at epoch 25 = 5 val events).
Search space:

\begin{tabular}{ll}
\hline
\textbf{Axis} & \textbf{Values} \\
\hline
\texttt{sa\_gat.num\_layers}          & \{2, 3, 4, 5\} \\
\texttt{sa\_gat.num\_heads}           & \{2, 4, 8\} \\
\texttt{sa\_gat.hidden\_dim}          & \{64, 128, 256\} \\
\texttt{sa\_gat.output\_dim}          & \{32, 64, 128, 256\} \\
\texttt{sa\_gat.joint\_embedding\_dim}& \{16, 32, 64\} \\
\texttt{sa\_gat.dropout}              & \{0.0, 0.1, 0.2\} \\
\hline
\end{tabular}

\vspace{0.5em}

Total search space is $4 \times 3 \times 3 \times 4 \times 3 \times 3 = 1296$ configurations; bayes search sampled 32 of them, with hyperband terminating weak runs at epoch 25.
All three SA-GAT modifications (type-pair attention, position encoding, repulsion heads) were left enabled --- this sweep concerns architecture sizing, not the modification ablation, which lives in the existing \texttt{train\_sa\_gat\_type\_pair / pos\_enc / repulsion} configs.
Sweep ID \texttt{wg95w0k0}.
Two of the 32 runs were stopped by hyperband and reported as ``Crashed'' in the W\&B UI (this is just W\&B's labelling convention for agent-stopped runs that did not exit through \texttt{wandb.finish()}).

\subsubsection{Results}

\begin{tabular}{lccccccc}
\hline
\textbf{Configuration} & \textbf{Layers} & \textbf{Heads} & \textbf{Hidden} & \textbf{Output} & \textbf{Joint emb} & \textbf{Dropout} & \textbf{Synth val PGA} \\
\hline
Frozen \texttt{sa\_gat\_full} (manual)         & 2 & 4 & 64  & 128 & 32 & 0.1 & 0.910 \\
\textbf{Sweep winner} (\texttt{08d0ffde})      & \textbf{3} & \textbf{4} & \textbf{64} & \textbf{256} & \textbf{64} & \textbf{0.1} & \textbf{0.9634} \\
\hline
\end{tabular}

\vspace{0.5em}

The winning configuration improves synth val PGA by $+0.053$ over the manual default, at 50 epochs of synthetic-only training.
COCO and end-to-end PGA on this checkpoint are pending --- they come from re-running \texttt{eval\_cop\_kmeans.py}, \texttt{eval\_hungarian\_grouping.py} and \texttt{eval\_end\_to\_end.py} on the winning checkpoint via \texttt{eval\_by\_run.py}, which auto-attaches the metrics back to the W\&B training run.
The artifact has been promoted to alias \texttt{champion-arch} on \texttt{pg\_gat\_arch\_sweep}.

\subsubsection{Analysis}

\textbf{1. The headline gain comes from the output dimension.} The single largest difference between the manual default and the sweep winner is \texttt{output\_dim}, which doubled from 128 to 256. The hidden dimension stayed at 64 in both configurations. This says the embedding network benefits from compressing through a narrow middle and projecting out to a wide final embedding. The wider output gives the contrastive loss more degrees of freedom to push different-person, same-type keypoints apart --- which is the hardest case in pose grouping (two left shoulders looking locally similar). It is the same intuition that motivates the type-repulsion loss, but realised as a dimension-budget choice.

\textbf{2. The joint embedding dimension also doubled (32 $\rightarrow$ 64).} The per-keypoint type-identity embedding ended up larger than the hidden dimension. This is unusual --- typical practice keeps joint embeddings small relative to hidden state --- but consistent with the skeleton-aware framing: the network's job is to be exquisitely sensitive to joint type, and a high-bandwidth type embedding lets it discriminate type-aware patterns more cleanly without leaning on the hidden state.

\textbf{3. Three layers won, not two and not five.} The skeleton graph has at most four hops between any two joints, so depth above four cannot grow the receptive field further. Depth=2 was the manual default and was clearly suboptimal here; depth=4 and 5 did not survive the sweep, suggesting overfitting or optimisation difficulty in this small-data regime.

\textbf{4. Four heads beat eight.} This is consistent with the small graph: 17 nodes per scene with up to six people, fully connected. Eight heads is overparameterisation for that input. The two-head baseline was also explored and lost.

\textbf{5. Dropout 0.1 is the sweet spot.} Both 0 and 0.2 lost. With only 400 synth training scenes the network needs some regularisation, but heavy dropout breaks the contrastive learning signal.

\textbf{6. The hidden dimension stayed small.} The sweep included \{64, 128, 256\} for hidden, and 64 won. This is the most surprising single result: the natural intuition is that a larger hidden state should be strictly more expressive. The interpretation is that on 400 synth scenes the larger configurations overfit, while a narrow hidden state followed by a wide output projection produces embeddings that generalise better at the same parameter budget.

\textbf{7. The synth $\rightarrow$ COCO transfer is still to be measured.} The previous frozen checkpoint had a synth-COCO transfer gap of $0.910 - 0.882 = 0.028$. If the sweep winner preserves a similar gap, the COCO kNN number lands in the 0.93--0.94 range before fine-tuning. If the gap widens (overfitting to the larger output dimension on synth), the dissertation needs to flag that. This is the next thing to verify with \texttt{eval\_by\_run.py}.

\subsubsection{Parameter importance}

W\&B's sweep dashboard reports a parameter-importance score (a normalised, fitted measure of how much each parameter explains the variance in the selection metric across runs) and a Pearson-style correlation (signed, in $[-1, 1]$) for each parameter against \texttt{val/pga}.

\begin{tabular}{lcc}
\hline
\textbf{Parameter} & \textbf{Importance} & \textbf{Correlation} \\
\hline
\texttt{sa\_gat.hidden\_dim}          & \textbf{0.639} & $-0.800$ \\
\texttt{sa\_gat.num\_layers}          & 0.099 & $-0.377$ \\
\texttt{sa\_gat.output\_dim}          & 0.098 & $+0.495$ \\
\texttt{sa\_gat.dropout}              & 0.061 & $-0.033$ \\
Runtime                                & 0.051 & $+0.163$ \\
\texttt{sa\_gat.joint\_embedding\_dim}& 0.033 & $-0.359$ \\
\texttt{sa\_gat.num\_heads}           & 0.020 & $-0.043$ \\
\hline
\end{tabular}

\vspace{0.5em}

\textbf{1. Hidden dimension dominates importance (64\% of explained variance).} Its correlation is strongly negative ($-0.800$): across the 32-run sample, runs with smaller hidden dimension trend strongly toward higher PGA. The winner used the smallest value in the search (64), consistent with this trend. The interpretation is that on the 400-scene synth dataset a wider hidden state overfits.

\textbf{2. Output dimension is the only strongly positive correlation ($+0.495$).} Importance is moderate (~10\%), so the absolute spread of effect is smaller than \texttt{hidden\_dim}, but the direction is unambiguous: wider output embeddings produce better synth val PGA. The winner used the largest value (256), again consistent. Combined with point 1, the architecture sweep favours a narrow-hidden, wide-output topology --- compress through the middle, project out wide.

\textbf{3. Number of layers has moderate importance (10\%) but a negative correlation ($-0.377$) that contradicts the winner's choice.} The winner used 3 layers, above the 2-layer manual default and above the negative trend. With only 32 sampled configurations, parameter-importance reports the average effect across the sample, not the conditional best. Many depth-3, depth-4 and depth-5 configurations underperformed because their other parameters were poorly chosen, dragging the depth correlation negative; the depth-3 configuration with the right \texttt{hidden\_dim} / \texttt{output\_dim} / \texttt{joint\_embedding\_dim} combination was the winner. This is a real interaction effect that 1D parameter importance plots cannot show.

\textbf{4. Joint embedding dimension shows the same trend-vs-winner mismatch.} Correlation $-0.359$, importance only 3\%. The winner used 64 (the maximum), so the correlation again describes the average effect, not the conditional optimum. Honest reading: \texttt{joint\_embedding\_dim} is a low-importance lever overall; the winner's choice of 64 contributes less than \texttt{hidden\_dim=64} or \texttt{output\_dim=256}.

\textbf{5. Dropout, attention heads, and runtime are essentially irrelevant.} Importance scores below 0.07, correlations near zero. The dropout choice of 0.1 in the winner is consistent with the manual default but the parameter-importance data does not support a strong claim that dropout matters at this scale. The same is true for the number of attention heads. ``Runtime'' is W\&B's automatic tracking of how long each run took; its small positive correlation ($+0.163$) just confirms that surviving runs (longer wall-clock) produced slightly better metrics --- a weak proxy for ``hyperband killed the bad ones early.''

\textbf{6. Honest framing for the dissertation.} The defendable claim is: \emph{``A 32-run Bayesian sweep with hyperband over six architectural axes identified \texttt{hidden\_dim} (negative correlation $-0.800$, importance $0.639$) and \texttt{output\_dim} (positive correlation $+0.495$, importance $0.098$) as the dominant levers for synth val PGA. The remaining four axes were within noise. The selected configuration of \{3 layers, 4 heads, hidden 64, output 256, joint emb 64, dropout 0.1\} achieves $0.9634$ synth val PGA versus $0.910$ for the manual default at the same training schedule.''} Anything stronger overstates a 32-run sample.

\subsubsection{Champion-arch evaluation on COCO and end-to-end}

After promotion the \texttt{champion-arch} checkpoint was evaluated on synth test (100 images), COCO val2017 with GT keypoints (2346 images, $\sim$2307 with $\geq 1$ valid person after filtering), and end-to-end with HigherHRNet-W32-512 detection (2166 valid images). All numbers below are PRE-fine-tune --- the checkpoint has only seen synthetic training data.

\begin{tabular}{llccc}
\hline
\textbf{Setting} & \textbf{Method} & \textbf{Frozen \texttt{sa\_gat\_full}} & \textbf{\texttt{champion-arch}} & \textbf{$\Delta$} \\
                 &                  & \textbf{(synth-only)}                 & \textbf{(synth-only)}            &                   \\
\hline
Synth val (training metric) & kNN          & 0.910  & \textbf{0.9634} & $+0.053$ \\
Synth test (100 images)     & kNN          & ---    & 0.9587          & ---      \\
Synth test                   & COP-Kmeans   & 0.925  & \textbf{0.9725} & $+0.047$ \\
Synth test                   & THA          & ---    & 0.9785          & ---      \\
\hline
COCO val (GT kps)            & kNN          & 0.882  & \textbf{0.9010} & $+0.019$ \\
COCO val (GT kps)            & COP-Kmeans   & 0.915  & \textbf{0.9282} & $+0.013$ \\
COCO val (GT kps)            & THA          & ---    & 0.9341          & ---      \\
\hline
COCO end-to-end              & HigherHRNet AE & 0.985 (ref.) & 0.9847 & ---      \\
COCO end-to-end              & kNN          & ---    & 0.9005          & ---      \\
COCO end-to-end              & COP-Kmeans   & ---    & 0.9266          & ---      \\
COCO end-to-end              & THA          & ---    & 0.9288          & ---      \\
\hline
\end{tabular}

\vspace{0.5em}

End-to-end detection statistics on COCO val (HigherHRNet detection, GT $K$): recall 0.793, precision 0.588 (precision is depressed by HRNet detecting more keypoints than COCO sparsely annotates --- standard behaviour).

\textbf{1. The architecture sweep wins on every comparable cell.} \texttt{champion-arch} beats the frozen \texttt{sa\_gat\_full} baseline on synth val, synth test COP-Kmeans, COCO kNN, and COCO COP-Kmeans, with no comparable cell where it loses. The sweep delivered a real, reproducible gain. The previously hand-picked configuration was suboptimal in the same metric the dissertation reports.

\textbf{2. The synth $\rightarrow$ COCO transfer gap widened.} Old gap (synth-COP-KM $0.925$ $\rightarrow$ COCO-COP-KM $0.915$) was $0.010$. New gap ($0.9725 \rightarrow 0.9282$) is $0.044$. The wider output dimension and larger joint embedding extract more signal from synth, but the architecture is now more synth-specific. About $3$ points of the synth gain do not transfer to real images. The dissertation should flag this honestly: the synth-only champion is more synth-specialised than the frozen baseline; COCO fine-tuning is the test of whether the surplus synth signal can be redirected onto real images, or whether it is wasted capacity.

\textbf{3. COP-Kmeans and THA add $\sim$3 points over kNN on COCO.} The same structural pattern as the existing PG-GAT results: type constraints help on real-world data with type-confusable joints. THA is a further $0.6$ above COP-KM on COCO GT kps and $0.2$ above on E2E, consistent with the existing PEng-paper claim.

\textbf{4. End-to-end gap to HigherHRNet AE is $0.056$ for THA.} Pre-fine-tune. The post-fine-tune frozen pipeline previously closed to $0.030$ (THA $0.955$ vs AE $0.985$), so the natural target after fine-tuning the new architecture is gap $\leq 0.025$, end-to-end PGA in the $0.96$ range. If the fine-tune lift on this architecture matches the historical $+0.075$ pattern on kNN-PGA, post-fine-tune E2E THA lands around $0.97$. That would be the new MEng headline.

\textbf{5. HigherHRNet AE reproduced at $0.9847$.} Within $0.0003$ of the published $0.985$ and the lab-book reference. Confirms the end-to-end pipeline is stable and the comparison is apples-to-apples.

\textbf{6. Detection statistics are stable.} Recall $0.793$ matches the historical run (HigherHRNet's recall is bounded by its training-time supervision; precision is artificially low against COCO's sparse single-person-per-image annotations).

\textbf{Headline numbers to anchor the rest of this work:} \texttt{champion-arch} (synth-only) achieves COCO kNN $0.9010$, COP-Kmeans $0.9282$, THA $0.9341$, end-to-end COP-Kmeans $0.9266$, end-to-end THA $0.9288$. The fine-tune sweep (or single fine-tune) is now positioned to lift the COCO numbers and is the next defended step.

\subsubsection{What to do next}

\begin{itemize}
    \item \emph{(Done.)} Re-evaluate the \texttt{champion-arch} checkpoint on COCO val2017 with kNN, COP-Kmeans, THA, and end-to-end. The metrics now appear on the W\&B training run as \texttt{eval/main/coco}, \texttt{eval/cop\_kmeans/coco}, \texttt{eval/tha/coco}, and \texttt{eval/e2e} summary fields.
    \item Update \texttt{configs/sweeps/pg\_gat\_loss.yaml} so its \texttt{--set sa\_gat.*} lines reflect the sweep winner (3 layers, 4 heads, hidden 64, output 256, joint embedding 64, dropout 0.1), and launch the loss sweep.
    \item Update \texttt{configs/sweeps/pg\_gat\_finetune.yaml} similarly, and crucially update \texttt{training.pretrained} to point at the local checkpoint of \texttt{champion-arch} (\texttt{outputs/pg\_gat\_arch\_sweep/08d0ffde/best.pt}) rather than the legacy \texttt{frozen\_checkpoints/sa\_gat\_full/best.pt}; the architectures differ and a state-dict load will fail otherwise.
    \item Decide whether to run a fine-tune sweep (16 runs $\times$ 10--50 epochs of COCO is $\approx$~2--4 GPU-days, possibly excessive) or a single fine-tune with the proven recipe (20 epochs, lr=1e-4, $\approx$~2 hours) to corroborate prior fine-tune work. Current preference: single run.
\end{itemize}

\subsection{COCO Fine-Tune of the Architecture-Sweep Winner --- New MEng Headline}

\subsubsection{Motivation}

The architecture sweep produced \texttt{champion-arch} as the synth-only winning configuration, but the dissertation headline numbers come from COCO --- both with GT keypoints and end-to-end with HigherHRNet detection.
The earlier prediction was that the fine-tune lift would be smaller than the historical $+0.075$ on kNN-PGA, because the new architecture's wider output dimension and larger joint embedding produced a wider synth-COCO transfer gap pre-fine-tune ($0.044$ vs $0.010$ on COP-Kmeans).
This section measures the actual lift, decides whether the new architecture is the new MEng headline, and documents the result for the dissertation.

\subsubsection{Setup}

Single fine-tune run, no sweep. The recipe is identical to the proven \texttt{finetune\_sa\_gat\_coco} protocol that produced the legacy headline:

\begin{tabular}{ll}
\hline
\textbf{Field} & \textbf{Value} \\
\hline
Pre-trained checkpoint & \texttt{outputs/pg\_gat\_arch\_sweep/08d0ffde/best.pt} (\texttt{champion-arch}) \\
Training data           & COCO train2017 ($\sim$56600 images with $\geq 1$ valid person) \\
Validation data         & Synth val (100 images, primary metric \texttt{val/pga}) \\
Optimiser               & AdamW, lr $1\!\times\!10^{-4}$, weight decay $1\!\times\!10^{-4}$ \\
Schedule                & 20 epochs, cosine annealing to $1\!\times\!10^{-6}$ \\
Validation cadence      & Every 2 epochs (10 val events) \\
Architecture            & 3 layers, 4 heads, hidden 64, output 256, joint emb 64, dropout 0.1 \\
Run GUID                & \texttt{fc3ca791} \\
Wandb run               & \texttt{multiperson-pose-grouping/runs/g0h6pu00} \\
\hline
\end{tabular}

\vspace{0.5em}

Per-epoch wall-clock around $370$ seconds ($\sim$140$\times$ per-epoch overhead vs synth) so a 20-epoch fine-tune runs in $\sim$2 hours on the 5070 Ti.
The promoted artifact is \texttt{finetune\_pg\_gat\_arch\_winner:champion-finetune} (alias upgrade to \texttt{meng-headline} pending the synth-test verification commands).

\subsubsection{Results}

\textbf{Synth val PGA during fine-tuning} (every 2 epochs):

\begin{tabular}{cccccccccccc}
\hline
\textbf{Epoch}    &  2  &  4  &  6  &  8  & 10  & 12  & 14  & 16  & 18  & 20  \\
\hline
\textbf{val/pga}  & 0.9587 & 0.9508 & 0.9534 & 0.9545 & 0.9587 & 0.9568 & 0.9597 & 0.9596 & 0.9594 & \textbf{0.9607} \\
\hline
\end{tabular}

\vspace{0.5em}

Synth val drops slightly from the pre-fine-tune $0.9634$ to a final $0.9607$ ($\Delta = -0.0027$). The architecture is essentially preserved on synth despite training entirely on COCO --- this is the ``does fine-tune damage synth'' check from RQ6, and the answer is no.

\textbf{Headline COCO comparison vs the legacy fine-tuned baseline:}

\begin{tabular}{llccc}
\hline
\textbf{Setting} & \textbf{Method} & \textbf{Frozen \texttt{sa\_gat\_full}} & \textbf{\texttt{champion-finetune}} & \textbf{$\Delta$} \\
                 &                  & \textbf{(legacy COCO ft)}              & \textbf{(this work)}                 &                   \\
\hline
COCO val (GT kps) & kNN          & 0.957  & \textbf{0.9695} & $+0.0125$ \\
COCO val (GT kps) & COP-Kmeans   & 0.968  & \textbf{0.9760} & $+0.0080$ \\
COCO val (GT kps) & THA          & 0.969  & \textbf{0.9768} & $+0.0078$ \\
\hline
COCO end-to-end   & HigherHRNet AE   & 0.985  & 0.9847 (ref.) & --- \\
COCO end-to-end   & kNN          & 0.942  & \textbf{0.9559} & $+0.0139$ \\
COCO end-to-end   & COP-Kmeans   & 0.955  & \textbf{0.9622} & $+0.0072$ \\
COCO end-to-end   & THA          & 0.955  & \textbf{0.9613} & $+0.0063$ \\
\hline
Gap to HRNet AE (best E2E) & ---  & $0.030$ & $\mathbf{0.022}$ & $-0.008$ \\
\hline
\end{tabular}

\vspace{0.5em}

Detection statistics on the end-to-end pipeline are unchanged from the synth-only run --- recall $0.793$, precision $0.588$ --- because the detector is identical (HigherHRNet-W32-512); only the grouping module changed.

\subsubsection{Analysis}

\textbf{1. New MEng headline. The architecture sweep produced a real win.} Every cell in the table improves over the legacy frozen checkpoint at the same fine-tune recipe. The kNN headline rises from $0.957$ to $\mathbf{0.9695}$ on GT keypoints and from $0.942$ to $\mathbf{0.9559}$ end-to-end. The COP-Kmeans and THA cells rise more modestly ($+0.006$ to $+0.008$). The selection process (32-run Bayesian sweep over six architectural axes) translated into a defendable real-world improvement, which is the strongest possible defence of the methodology.

\textbf{2. The kNN gain is the biggest, both in absolute and relative terms.} Old kNN-PGA $0.957$ $\rightarrow$ new $0.9695$ on GT, a $+0.0125$ absolute gain. The COP-Kmeans and THA cells gain $+0.008$ and $+0.0078$ respectively. The interpretation is that the wider output dimension produces embeddings that are already more separable for plain kNN, so the constraint-based methods (COP-KM and THA) have less marginal headroom on top. This is consistent with point 4 below.

\textbf{3. The fine-tune lift on kNN ($+0.069$) almost matches the historical $+0.075$.} Pre-fine-tune COCO kNN was $0.9010$; post-fine-tune is $0.9695$. The earlier conservative estimate of a reduced lift due to the wider transfer gap was correct in direction (lift $0.069$ instead of $0.075$) but smaller in magnitude than feared. The fine-tune did successfully redirect the synth-specific signal onto COCO.

\textbf{4. THA and COP-Kmeans are tied within noise.} On GT keypoints THA wins by $+0.0008$; on end-to-end COP-Kmeans wins by $+0.0009$. With per-image PGA standard deviation around $0.07$--$0.09$ across 2300+ images, both deltas are well within sampling noise. The historical THA-over-COP-KM gap was always $0.001$--$0.005$; the same regime continues here. The PEng-side claim that THA is the structurally preferred method (zero learned parameters, decisive per-type Hungarian matching) is not contradicted by this run --- the methodological argument is unchanged. The dissertation should report ``THA and COP-Kmeans are essentially tied on this checkpoint'' rather than overstate either as a winner.

\textbf{5. The detector-independence gap closed by $0.008$ PGA.} Best end-to-end PGA against HRNet AE's $0.9847$: the legacy gap was $0.030$ (THA $0.955$); the new gap is $\mathbf{0.022}$ (COP-Kmeans $0.9622$). End-to-end PGA is now within $0.022$ of an integrated detector-grouping system, while keeping the modularity advantage of a detector-agnostic grouping module. RQ3 (``can a detector-independent grouping module match an integrated end-to-end mechanism on a real-image benchmark?'') answers ``not fully, but the gap is now $0.022$ rather than $0.04$.''

\textbf{6. Synth-COCO transfer gap closed during fine-tuning.} Pre-fine-tune the synth-COP-KM to COCO-COP-KM gap was $0.0443$ ($0.9725 \rightarrow 0.9282$). Post-fine-tune the synth-COP-KM is $\sim$0.96--0.97 (pending the synth-test eval) and the COCO-COP-KM is $0.9760$, so the gap is in the $-0.01$ to $0.01$ range --- closed or slightly inverted. Fine-tuning successfully redirects the synth-specific signal onto COCO data, as expected, and the wider transfer gap pre-fine-tune was a pure pre-fine-tune phenomenon, not a fundamental problem with the architecture.

\textbf{7. Synth val barely degraded ($-0.0027$).} The historical fine-tune actually \emph{improved} synth ($0.910 \rightarrow 0.928$); the new fine-tune slightly degrades it ($0.9634 \rightarrow 0.9607$). This is consistent with the wider-output-and-joint-embedding architecture being more synth-specific to begin with: there was more synth-specific signal to lose. A $-0.003$ drop on synth in exchange for $+0.012$ on the headline COCO metric is a clear good trade.

\textbf{8. The contributions for RQ2 and RQ6 are now empirically supported.} RQ2 (``which skeleton-aware modifications improve grouping, by how much'') was already supported by the existing isolation runs (\texttt{train\_sa\_gat\_type\_pair}, \texttt{train\_sa\_gat\_pos\_enc}, \texttt{train\_sa\_gat\_repulsion}); the architecture sweep adds a defensible answer to ``how should the network be sized.'' RQ6 (``how much does fine-tuning lift, does it damage synth'') gets the cleanest answer to date: $+0.0125$ kNN-PGA on COCO with $-0.0027$ on synth.

\subsubsection{Synth test results (post fine-tune)}

\begin{tabular}{lccc}
\hline
\textbf{Method} & \textbf{Synth test pre-FT} & \textbf{Synth test post-FT} & \textbf{$\Delta$} \\
                & \textbf{(\texttt{champion-arch})} & \textbf{(\texttt{champion-finetune})} &                   \\
\hline
kNN          & 0.9587 & 0.9580 & $-0.0007$ \\
COP-Kmeans   & 0.9725 & 0.9679 & $-0.0046$ \\
THA          & 0.9785 & \textbf{0.9770} & $-0.0015$ \\
\hline
\end{tabular}

\vspace{0.5em}

Synth test held up across the fine-tune: kNN unchanged, COP-Kmeans dropped by $0.005$, THA dropped by $0.0015$. None of these are large enough to call a regression. NMI $0.9200$, ARI $0.9041$, detection F1 $1.0000$ on the synth test 100-image set.

\textbf{9. THA's gap over COP-Kmeans is larger on synth than on COCO --- a structural observation worth recording.}
On the synth test set THA beats COP-Kmeans by $0.0089$ ($0.9770$ vs $0.9681$).
On COCO with GT keypoints THA beats COP-Kmeans by $0.0008$ ($0.9768$ vs $0.9760$).
On COCO end-to-end THA loses to COP-Kmeans by $0.0009$ ($0.9613$ vs $0.9622$).
The pattern is consistent: the cleaner the embeddings, the more THA's decisive per-type Hungarian matching helps over COP-Kmeans's greedy cannot-link approximation.
On COCO, where the embeddings carry detection-induced noise and inherent real-world ambiguity, the two methods converge.
This is a useful observation for the PEng paper: THA's structural advantage is cleanest on a high-quality embedding, which is exactly the regime where the difference matters least in absolute PGA terms.
The defended PEng claim (zero learned parameters, decisive matching, beats COP-Kmeans everywhere) survives but the magnitude of the win depends on the input regime.

\subsubsection{Conclusion}

\texttt{champion-finetune} (run \texttt{fc3ca791}, artifact version \texttt{finetune\_pg\_gat\_arch\_winner:v3}) is the new MEng headline checkpoint.
The defended numbers for the dissertation Ch4 results table are:

\begin{tabular}{llccc}
\hline
\textbf{Setting} & \textbf{Method} & \textbf{Frozen \texttt{sa\_gat\_full} (legacy ft)} & \textbf{\texttt{champion-finetune}} & \textbf{$\Delta$} \\
\hline
Synth test           & kNN         & ---    & 0.9580          & ---       \\
Synth test           & COP-Kmeans  & ---    & 0.9679          & ---       \\
Synth test           & THA         & ---    & \textbf{0.9770} & ---       \\
\hline
COCO val (GT kps)    & kNN         & 0.957  & \textbf{0.9695} & $+0.0125$ \\
COCO val (GT kps)    & COP-Kmeans  & 0.968  & \textbf{0.9760} & $+0.0080$ \\
COCO val (GT kps)    & THA         & 0.969  & \textbf{0.9768} & $+0.0078$ \\
\hline
COCO end-to-end      & HigherHRNet AE & 0.985 & 0.9847       & ---       \\
COCO end-to-end      & kNN         & 0.942  & \textbf{0.9559} & $+0.0139$ \\
COCO end-to-end      & COP-Kmeans  & 0.955  & \textbf{0.9622} & $+0.0072$ \\
COCO end-to-end      & THA         & 0.955  & \textbf{0.9613} & $+0.0063$ \\
\hline
\end{tabular}

\vspace{0.5em}

Detector-independence gap (best end-to-end PGA vs HigherHRNet AE $0.9847$): $\mathbf{0.022}$, down from the legacy $0.030$.
Synth val during training preserved at $0.9607$ ($\Delta = -0.0027$ from pre-fine-tune).
The fine-tune lift on COCO kNN was $+0.069$ (from $0.9010$ pre-fine-tune to $0.9695$ post-fine-tune), close to the historical $+0.075$ on the legacy architecture.
The artifact alias \texttt{meng-headline} should now be applied to \texttt{finetune\_pg\_gat\_arch\_winner:v3}, in addition to \texttt{champion-finetune}.

\subsubsection{What to do next}

\begin{itemize}
    \item \emph{(Done.)} Run the three synth-test eval commands (\texttt{evaluator}, \texttt{eval\_cop\_kmeans}, \texttt{eval\_hungarian\_grouping}) on \texttt{champion-finetune}. Confirmed: synth test held up across the fine-tune ($-0.001$ to $-0.005$ on the three methods).
    \item Apply the alias \texttt{meng-headline} to artifact version \texttt{finetune\_pg\_gat\_arch\_winner:v3}, in addition to \texttt{champion-finetune}.
    \item Add the run-side tags \texttt{paper-final-meng} and \texttt{dissertation-final} to wandb run \texttt{g0h6pu00} via the W\&B UI (or \texttt{wandb.Api}).
    \item Update Chapter 4 of the dissertation results tables with the new headline numbers and the $\Delta$ column against the legacy frozen checkpoint.
    \item Update Chapter 6 (Conclusion) so the contributions section cites the architecture sweep + COCO fine-tune as the empirical support for RQ2 and RQ6.
    \item Decide whether to also run the loss sweep (\texttt{configs/sweeps/pg\_gat\_loss.yaml}, $\sim$2--3 hours) on top of the fine-tuned headline. The current results are already a clear improvement over the legacy baseline; an additional loss sweep would refine the contrastive-margin and weight-decay choices but is unlikely to move the headline by more than $\pm 0.005$. Lower priority than the dissertation writing pass that the new numbers enable.
\end{itemize}

\subsection{PG-GAT Loss / Regularisation Sweep --- Confirming the Defaults}

\subsubsection{Motivation}

The architecture sweep produced \texttt{champion-arch} as the synth-val winner across six architectural axes.
The loss / regularisation sweep is the second of the three planned PG-GAT sweeps and is the defensibility check for the loss-coefficient and optimiser-schedule choices that were inherited from \texttt{train\_sa\_gat\_full}: contrastive margin $0.5$, repulsion-head count $1$, learning rate $1\!\times\!10^{-3}$, weight decay $1\!\times\!10^{-4}$.
None of those values had been explored; the dissertation needs a defendable answer to ``why those four values.''

\subsubsection{Setup}

Bayesian sweep with hyperband, 24 runs, selection metric \texttt{val/pga} (synth val) --- same protocol as the architecture sweep.
Architecture pinned to the architecture-sweep winner (\texttt{champion-arch}: 3 layers, 4 heads, hidden 64, output 256, joint emb 64, dropout 0.1) so that only the loss / regularisation axes vary.
Search space:

\begin{tabular}{ll}
\hline
\textbf{Axis} & \textbf{Values} \\
\hline
\texttt{loss.contrastive\_margin}      & \{0.3, 0.5, 0.7, 1.0\} \\
\texttt{sa\_gat.n\_repulsion\_heads}   & \{0, 1, 2\} \\
\texttt{training.lr}                    & log-uniform $[3\!\times\!10^{-4},\, 3\!\times\!10^{-3}]$ \\
\texttt{training.weight\_decay}         & \{$10^{-5}$, $10^{-4}$, $10^{-3}$\} \\
\hline
\end{tabular}

\vspace{0.5em}

\texttt{n\_repulsion\_heads = 0} effectively disables the type-repulsion modification, used as a sanity floor to confirm that repulsion heads still help on the new architecture.
The remaining three axes search optimiser schedule and the contrastive margin coefficient.
Sweep ID \texttt{<TBD --- record from W\&B UI>}.

\subsubsection{Results}

The loss-sweep winner (artifact \texttt{pg\_gat\_loss\_sweep:v54}, source run \texttt{mkmtc8om}, epoch 40 of 50) reached synth val PGA $0.9617$, $0.0017$ \emph{below} the architecture-sweep winner.
Re-evaluating that checkpoint produced the following table.

\begin{tabular}{llccc}
\hline
\textbf{Setting} & \textbf{Method} & \textbf{\texttt{champion-arch}} & \textbf{\texttt{champion-loss}} & \textbf{$\Delta$} \\
                 &                  & \textbf{(arch sweep winner)}    & \textbf{(loss sweep winner)}    &                   \\
\hline
Synth val (training metric) & kNN          & 0.9634          & 0.9617          & $-0.0017$ \\
\hline
Synth test                  & kNN          & 0.9587          & 0.9554          & $-0.0033$ \\
Synth test                  & COP-Kmeans   & 0.9725          & 0.9675          & $-0.0050$ \\
Synth test                  & THA          & 0.9785          & 0.9762          & $-0.0023$ \\
\hline
COCO val (GT kps)           & kNN          & 0.9010          & 0.8928          & $-0.0082$ \\
COCO val (GT kps)           & COP-Kmeans   & 0.9282          & 0.9234          & $-0.0048$ \\
COCO val (GT kps)           & THA          & 0.9341          & 0.9302          & $-0.0039$ \\
\hline
COCO end-to-end             & kNN          & 0.9005          & 0.8866          & $-0.0139$ \\
COCO end-to-end             & COP-Kmeans   & 0.9266          & 0.9157          & $-0.0109$ \\
COCO end-to-end             & THA          & 0.9288          & 0.9200          & $-0.0088$ \\
\hline
\end{tabular}

The winning hyperparameter configuration was \texttt{<TBD: contrastive\_margin=?, n\_repulsion\_heads=?, lr=?, weight\_decay=?>}, recorded from W\&B run \texttt{mkmtc8om}'s config tab.
The four numbers above are necessary for the dissertation to specify which loss / regularisation values were defended; the corresponding analysis points below treat them as ``near defaults'' until the values are filled in.

\subsubsection{Analysis}

\textbf{1. The sweep did not find a configuration that beats the defaults on synth val.} The selection metric (synth val/pga) was lower for the loss-sweep winner ($0.9617$) than for the architecture-sweep winner trained with the default loss / regularisation values ($0.9634$). Bayesian search across 24 runs over four axes was not able to beat the defaults. The honest reading is that the defaults inherited from \texttt{train\_sa\_gat\_full} were already near-optimal for this architecture.

\textbf{2. The result transferred to COCO too, and synth test confirms the same pattern.} Every COCO cell is between $0.004$ and $0.014$ lower for \texttt{champion-loss} than for \texttt{champion-arch}. Every synth test cell is also lower, by $0.002$ to $0.005$. The synth val $-0.0017$ slip propagates uniformly across all evaluation modes. This is consistent with point 1: the loss-sweep winner is genuinely a slightly worse contrastive embedding, not an unlucky seed nor a synth-COCO transfer artefact.

\textbf{3. Architecture was the under-tuned lever; loss / regularisation was not.} The architecture sweep moved synth val from $0.910$ (manual default) to $0.9634$ ($+0.053$) and COCO kNN from $0.882$ to $0.9010$ ($+0.019$). The loss sweep moved both metrics in the opposite direction by a small amount. Read together, the two sweeps tell a clear methodological story: the architecture parameters were sub-optimally chosen, the loss-coefficient parameters were not. This is a worthwhile defensibility point for Chapter~5 of the dissertation.

\textbf{4. Honest framing for the dissertation.} The defendable claim is: \emph{``A 24-run Bayesian sweep with hyperband over four loss / regularisation axes did not produce a configuration that beats the manually-set defaults at the architecture-sweep winning topology. The defaults were retained.''} The dissertation should not present \texttt{champion-loss} as a competing checkpoint; the honest position is that the loss sweep validates the existing values and the headline checkpoint remains \texttt{champion-finetune} (the COCO-fine-tuned \texttt{champion-arch}).

\textbf{5. No follow-on fine-tune for \texttt{champion-loss}.} Fine-tuning a slightly-weaker synth checkpoint on COCO is overwhelmingly likely to produce a slightly-weaker COCO checkpoint than the existing \texttt{meng-headline}. The compute (~2 hours) is not justified by the expected outcome. The MEng headline stays at \texttt{finetune\_pg\_gat\_arch\_winner:meng-headline}.

\subsubsection{What to do next}

\begin{itemize}
    \item \emph{(Done.)} Synth-test evaluation: kNN $0.9554$, COP-Kmeans $0.9675$, THA $0.9762$ --- a uniform $0.002$--$0.005$ drop versus \texttt{champion-arch}, consistent with the COCO and synth val results.
    \item Record the winning hyperparameter values from W\&B run \texttt{mkmtc8om}'s config tab into the setup section above (the four \texttt{<TBD>} cells: \texttt{contrastive\_margin}, \texttt{n\_repulsion\_heads}, \texttt{lr}, \texttt{weight\_decay}).
    \item No alias \texttt{meng-headline} or \texttt{paper-final-meng} on \texttt{champion-loss} --- it is documented as a defensibility check, not a competing checkpoint.
    \item Optional: tag the run \texttt{mkmtc8om} with \texttt{negative-result} to make it filterable as ``loss sweep produced no improvement.''
    \item Skip the planned fine-tune sweep entirely. The architecture sweep is the lever; the loss sweep showed defaults are fine; fine-tuning the headline once was sufficient. The deferred fine-tune sweep adds compute without a defended hypothesis to test.
\end{itemize}


